\section*{v1.7.0 Closure--Projection Synthesis Core}
\addcontentsline{toc}{section}{v1.7.0 Closure--Projection Synthesis Core}

\begin{remark}[Normalization note for v1.7.0]
\label{rem:normalization-note-v170}
In the present release, the mathematical content of this synthesis layer is subsumed by the OP transport theorem and the OP3 residual/shell interface. The block is retained as historical compression logic and interpretive guidance. This synthesis block is therefore kept as a historical transport precursor: its closure--projection synthesis is now carried proof-theoretically by Phase~II-D, Phase~II-E, and Phase~II-F, and here it should be read as an earlier integrative summary rather than as an independent axiomatic center.
\end{remark}

This release adds a further synthesis layer between the closure blocks already
introduced and the inherited structural core of the monolith. Its purpose is to
state, in one explicit place, how closure statements are transported across the
probe sector, how projection claims inherit their admissibility status from the
closure architecture, and how appendix tables are to be used as reversible proof
interfaces rather than as detached summaries.

\subsection*{1. Closure-to-projection transport discipline}
The central operational rule of v1.7.0 is that no projection statement may be
read as primary if its admissibility has not first been transported through the
closure chain. In other words, the order of explanation remains
\[
\begin{aligned}
&\text{local structure}
\Longrightarrow \text{HC}
\Longrightarrow \text{closure}\\[0.15em]
&\Longrightarrow \text{probe-compatible reduction}
\Longrightarrow \text{projection readout}.
\end{aligned}
\]
A projection is therefore not an independent object but a compressed exposure of
closure-stable content under a chosen observational or effective frame. This
rule keeps the paper aligned with the structure-first programme and prevents
later phenomenological language from silently replacing the closure mechanism
that generated it.

\paragraph{Synthesis Proposition A (closure status survives transport).}
If a statement has been established in the closure regime and is moved to a
probe-compatible or projective representation by a transport that preserves the
common body, then its structural status is not lost. What changes is the mode of
appearance, not the underlying admissibility class.

\paragraph{Proof sketch.}
The common body is fixed by the ground condition and stabilized by HC. Any
transport that preserves this body may change coordinates, phase labels, or
sector language, but it does not create a new ontology. Hence closure status is
inherited through the transport.

\subsection*{2. Probe reduction as a controlled interface}
The probe sector is not only a search device; in v1.7.0 it is also an
interface theorem. Two probes moving against one another, each carrying triolic
organization and coupled to an orthogonal third direction, define a reduced test
geometry in which admissible transitions can be filtered before full projection.
This means that the probe sector is a disciplined compression stage: it narrows
candidate structure while still remaining within closure control.

\paragraph{Synthesis Proposition B (probe reduction preserves closure discipline).}
Whenever the probe sector is constructed from triolic substructures under the
standing orthogonality and HC rules, the reduced search geometry cannot produce
admissible outputs that violate closure. It may discard candidates, rank them,
or expose only a subset of modes, but it cannot certify a structurally illegal
sector.

\paragraph{Proof sketch.}
The probe geometry is not external to the closure architecture. It is built from
substructures already governed by HC and orthogonality. Therefore any reduction
performed inside the probe channel is a reduction of representation, not a
release from confinement.

\subsection*{3. Projection tables as reversible proof interfaces}
The appendix tables introduced in the previous versions are to be read neither
as cosmetic summaries nor as final stand-alone derivations. Their exact role is
reversible: each table entry points upward to the closure and lemma chain from
which it is licensed, and downward to the sector language in which it may be
used. In this sense, a table row functions as a compact proof interface.

\paragraph{Synthesis Principle (table row admissibility).}
A table row is admissible only if one can move in both directions:
\begin{enumerate}[leftmargin=2em,label=(\alph*)]
  \item upward, from the row to the closure/lemma chain that stabilizes it; and
  \item downward, from the row to a projective or phenomenological reading that
  does not exceed the transported status of the original statement.
\end{enumerate}
If either direction is unavailable, the row is only heuristic and must be marked
as such.

\subsection*{4. Micro-to-macro synthesis route}
The micro-to-macro route is now stated in a more disciplined form. The local
triolic unit is not yet a macroscopic law, but it already fixes the admissible
coupling pattern. HC then stabilizes multiple appearances of the same body,
closure confines the admissible dynamics, the probe sector performs controlled
reduction, and only then do effective macroscopic equations appear as readouts.
The macro level is therefore a closure-saturated exposure of the micro level,
not a rival explanatory layer.

\paragraph{Synthesis Proposition C (macro equations are closure-saturated readouts).}
Under the standing assumptions of the monolith, any effective macroscopic
relation used later in the paper must be readable as a closure-saturated image
of the local structural chain. Hence macro equations are legitimate exactly to
the degree that the closure-to-projection transport remains traceable.

\subsection*{5. Appendix back-reference discipline, strengthened}
The appendix is now explicitly promoted to a back-reference layer. This means
that every compressed proof move used in the running text should be recoverable
through one of three routes: a lemma reference, a tool reference, or a table
reference with reversible upward justification. The appendix is therefore not an
archive placed after the fact; it is the controlled external memory of the proof
architecture.

\paragraph{Operational rule for later versions.}
Whenever a derivation step is abbreviated in the main line, the abbreviated step
should be tagged by a reusable interface name that can be expanded again from
the appendix without changing the logical order of the argument. This preserves
speed without sacrificing structural auditability.


\section*{v1.8.0 Structural Interface Consolidation Core}
\addcontentsline{toc}{section}{v1.8.0 Structural Interface Consolidation Core}

\begin{remark}[Normalization note for v1.8.0]
\label{rem:normalization-note-v180}
This interface-consolidation block is retained as historical interface memory. Its active proof-bearing content is now normalized into the semantic/tensor interface and grouped transport phases, so the present block should be read as supporting consolidation rather than as a competing interface foundation.
\end{remark}

This release adds a further consolidation layer between the inherited closure--projection
synthesis core and the longer structural body of the monolith. The aim of
v1.8.0 is to make the interfaces between closure statements, probe reductions,
projection tables, and appendix recoverability even more explicit, so that later
mathematical deepening can reuse the same transport rules without reopening the
basic admissibility logic.

\subsection*{1. Interface invariance principle}
The first strengthening of v1.8.0 is that every admissible transition between
main text, probe representation, and appendix table must preserve the same
structural body. This means that the interface is not merely notational. It is a
controlled transfer channel whose legitimacy depends on invariance of the common
closure-supported content.

\paragraph{Consolidation Proposition A (interface invariance).}
If two representations are linked by a transport that preserves the closure-fixed
body, then they are distinct expositions of the same admissible statement rather
than competing structural claims.

\paragraph{Proof sketch.}
The ground condition and HC fix the admissible body before any later compression.
A transport that keeps this body intact may change display, resolution, or
sector language, but it cannot alter the admissibility source. Hence interface
moves preserve statement identity at the structural level.

\subsection*{2. Probe filtration as a pre-projective gate}
The probe sector is now stated even more sharply as a pre-projective gate. Its
job is not to invent physical outputs but to filter, compress, and expose only
those modes that remain compatible with the closure discipline already fixed in
the local structural chain.

\paragraph{Consolidation Proposition B (probe filtration is gatekeeping, not generation).}
Under the standing orthogonality and HC assumptions, the probe sector may rank or
suppress candidate modes, but it cannot generate a physically admissible mode
whose closure ancestry is absent.

\paragraph{Proof sketch.}
The probes inherit their admissibility from the same triolic and closure-governed
body as the main sector. Therefore filtration can only reveal or reject what is
already licensed; it cannot certify an unlicensed sector by reduction alone.

\subsection*{3. Table-to-lemma reversibility}
The tables and tool blocks gathered in the appendix are promoted once more from
summary devices to reversible proof interfaces. For v1.8.0 this reversibility
is sharpened into a two-way requirement between rows, lemmas, and transport
principles.

\paragraph{Consolidation Principle (row/lemma reversibility).}
A compact row or tool entry is fully admissible only if the reader can reconstruct
both
\begin{enumerate}[leftmargin=2em,label=(\alph*)]
  \item the lemma or closure chain that licenses the row from above, and
  \item the bounded projection or interpretation that uses the row from below.
\end{enumerate}
Where only one direction is available, the entry remains heuristic and must be
marked as such.

\subsection*{4. Strengthened micro-to-macro readout discipline}
The micro-to-macro route is no longer only a synthesis slogan. In v1.8.0 it is
a regulated readout pipeline:
\[
\begin{aligned}
&\text{local triolic body} \Longrightarrow \text{HC stabilization} \Longrightarrow \text{closure confinement}\\
&\Longrightarrow \text{probe filtration} \Longrightarrow \text{table/interface compression} \Longrightarrow \text{effective projection}.
\end{aligned}
\]
This ordering is now part of the operational discipline of the document. A macro
statement is acceptable only to the extent that its position in this chain can be
traced backward without gap.

\paragraph{Consolidation Proposition C (macro readout requires backward traceability).}
Any effective large-scale law cited later in the monolith is structurally valid only
when it admits a backward trace through table/interface compression, probe
filtration, closure confinement, and HC stabilization to the local triolic body.

\subsection*{5. Appendix memory discipline, strengthened again}
The appendix is therefore not only a storage area for auxiliary derivations. It is
a controlled external memory for the proof architecture. In v1.8.0 every major
abbreviation introduced in the running text should be recoverable by a named
interface path: lemma path, tool path, table path, or transport path. This rule
prepares the next versions for denser mathematical unfolding without losing the
structure-first audit trail.

\paragraph{Operational rule for later versions.}
When a proof step is compressed in the main line, the compression should carry a
stable interface name whose recovery route is unique up to equivalent structural
transport. This keeps the monolith fast in presentation while remaining expandable
under later formal pressure.

\section{Structural Core and Closure Principles}
\label{sec:structural-core}
%% ============================================================

Before compressing the preceding results into a single ToE master equation, we
make explicit the structural kernel that has been implicit throughout the
construction.

\begin{definition}[Structural core]
\label{def:structural-core}
The structural core of the theory consists of the following coupled
requirements:
\begin{enumerate}
\item a common underlying field object $\Psi$;
\item a triolic decomposition
\[
  \Psi = \Psi_1\oplus\Psi_2\oplus\Psi_3;
\]
\item an orthogonal closure/light sector $\Psi_3$ relative to the first two
sectors;
\item a dual unitary symmetry exchanging $\Psi_1$ and $\Psi_2$;
\item a Heisenberg compensator $C$ emerging from the ground condition and
selecting the collapse/light output;
\item a projected closure condition determining the physically admissible
sector;
\item a phase-shifted observer frame in which effective quantities are read
out.
\end{enumerate}
\end{definition}

\begin{theorem}[Existence of a nontrivial closure-stable triolic system]
\label{thm:structural-existence}
Assume the structural core of Definition~\ref{def:structural-core}. Then there
exists a nontrivial class of configurations in which the triolic decomposition,
compensator consistency, and projected closure condition are simultaneously
satisfied.
\end{theorem}

\begin{proof}
The preceding sections already exhibit the required ingredients separately.
The central theorem provides a common ground object together with its triolic
unfolding and Lorentz-compatible strata. The unitary symmetry $J$ links the
first two sectors, while the Heisenberg compensator produces the
collapse/light sector as a structurally distinguished output. The subsequent
probe--closure analysis introduces the projected closure residue and identifies
an admissible closed sector by the condition that leakage into the complement
is suppressed. Since these ingredients are realized within one and the same
construction, the intersection of the triolic, compensator-consistent, and
closure-stable constraints is nonempty. Hence a nontrivial closure-stable
triolic system exists.
\end{proof}

\begin{proposition}[Centrality of the Heisenberg compensator]
\label{prop:hc-centrality}
Within the structural core, the Heisenberg compensator is not an auxiliary
correction term but the central structural hinge linking duality, collapse,
closure stability, and the phase-shifted observer description.
\end{proposition}

\begin{proof}
Without the duality symmetry, the first two sectors would remain merely
parallel descriptions. Without the compensator, their common collapse/light
output would not be selected. Without this selected output, the projected
closure condition would not isolate a physically admissible sector. And
without that closure-stable sector, the observer-frame description would lack a
well-defined quantity to read out after phase shift. Thus the compensator is
precisely the operator through which duality becomes collapse, collapse becomes
closure, and closure becomes observable effective physics.
\end{proof}


\subsection*{Preliminary Semantic and Tensorial Kernel Interface}
\addcontentsline{toc}{subsection}{Preliminary Semantic and Tensorial Kernel Interface}

This subsection states the semantic/tensor interface in anticipatory form. Its formal unfolding as a lemma--theorem chain is given later in Phase II-E.

The structural kernel can be stated on two coordinated levels. The first is a
semantic layer of sets, states, maps, and admissibility predicates. The second
is a tensorial realization layer in which the same constraints are written in a
notation compatible with the Riemann/Sphaera core. The tensor layer does not
replace the semantic one; it compresses it into a transportable operator form.

\begin{definition}[Semantic structural kernel]
\label{def:semantic-structural-kernel}
The semantic structural kernel consists of the data
\[
(\mathbb S,\{S_n\}_{n\ge 0},S_0,\hat K,J,C,\mathcal M)
\]
together with the admissibility predicates
\[
\begin{aligned}
&\mathrm{Triolic}(x),\ \mathrm{Anchored}(x),\ \mathrm{FrameShifted}(x),\ \mathrm{GroundLinked}(x),\\
&\mathrm{HCForced}(x),\ \mathrm{ClosureStable}(x),\ \mathrm{HigherConfined}(x),\ \mathrm{SeptinionReadable}(x).
\end{aligned}
\]
These objects satisfy the following semantic conditions:
\begin{enumerate}[leftmargin=2em,label=(S\arabic*)]
\item \textbf{Cascade hierarchy:}
\[
\mathbb S=S_N\supseteq S_{N-1}\supseteq \cdots \supseteq S_1\supseteq S_0,
\qquad
\hat K=\Pi_0\circ\Pi_1\circ\cdots\circ\Pi_{N-1}.
\]
\item \textbf{Ground-state terminality:}
\[
x\in S_0 \Longleftrightarrow \hat K(x)=x.
\]
\item \textbf{Maximal-totality closure rule:} if $M^*\cap \mathbb S=\emptyset$, then
\[
\hat K(M^*)\subseteq S_0.
\]
\item \textbf{$J$-involution:}
\[
J^2=\mathrm{id}.
\]
\item \textbf{Compensator map:}
\[
C(x):=\tfrac12\bigl(x+J(x)\bigr).
\]
\item \textbf{Triolic minimality:} if $\mathrm{Triolic}(x)$, then $x$ admits a three-sector
organization $(K_1,K_2,K_3)$.
\item \textbf{Orthogonal anchor:} if $\mathrm{Anchored}(x)$, then the third sector is
orthogonal to the first two,
\[
K_3\perp K_1,
\qquad
K_3\perp K_2.
\]
\item \textbf{Frame-shift principle:} $\mathrm{FrameShifted}(x)$ means that $x$ is a shifted
or regime-dependent readout of a common structural body rather than an
ontologically independent object.
\end{enumerate}
\end{definition}

\begin{definition}[Tensor realization of the structural kernel]
\label{def:tensor-structural-kernel}
A tensor realization of Definition~\ref{def:semantic-structural-kernel}
is given by the operator family
\[
\mathcal J^\mu{}_{\nu},\qquad
\mathcal C^\mu{}_{\nu},\qquad
\Pi_n{}^\mu{}_{\nu},\qquad
\hat K^\mu{}_{\nu},\qquad
\mathcal M^\mu{}_{\nu},
\]
subject to the following identities:
\begin{enumerate}[label=(T\arabic*)]
  \item \textbf{Involution:}
    $\mathcal J^\mu{}_{\alpha}\mathcal J^\alpha{}_{\nu}=\delta^\mu{}_{\nu}$.
  \item \textbf{Compensator idempotence:}
    $\mathcal C^\mu{}_{\alpha}\mathcal C^\alpha{}_{\nu}=\mathcal C^\mu{}_{\nu}$.
  \item \textbf{Cascade composition:}
    $\hat K^\mu{}_{\nu}=(\Pi_0\Pi_1\cdots\Pi_{N-1})^\mu{}_{\nu}$.
  \item \textbf{$\mathcal M$-admissibility (intertwining):}
    The Millennium matrix $\mathcal M^\mu{}_{\nu}$ is required to be
    \emph{admissible} in the sense that it preserves the algebraic
    structure of the kernel:
    \[
      \mathcal M^\mu{}_{\alpha}\mathcal J^\alpha{}_{\nu}
      =\mathcal J^\mu{}_{\alpha}\mathcal M^\alpha{}_{\nu},
      \qquad
      \mathcal M^\mu{}_{\alpha}\mathcal C^\alpha{}_{\nu}
      =\mathcal C^\mu{}_{\alpha}\mathcal M^\alpha{}_{\nu}.
    \]
    Whenever layer transport is required to preserve the cascade reading,
    additionally
    \[
      \mathcal M^\mu{}_{\alpha}\hat K^\alpha{}_{\nu}
      =\hat K^\mu{}_{\alpha}\mathcal M^\alpha{}_{\nu}.
    \]
\end{enumerate}
The tensor $\mathcal J$ represents the involutive dual/frame symmetry,
$\mathcal C$ represents the compensator projection, the $\Pi_n$ represent
cascade projections between adjacent strata, and $\mathcal M$ is the
Millennium matrix as an admissible structure-preserving converter between
admissible layers. Condition~(T4) is the defining admissibility condition
on $\mathcal M$; it is not derivable from (T1)--(T3) alone, but is imposed
as the minimal requirement for $\mathcal M$ to qualify as a
structure-preserving converter in the sense of
Definition~\ref{def:semantic-structural-kernel}.
\end{definition}

\begin{proposition}[Semantic forcing chain]
\label{prop:semantic-forcing-chain}
Within the semantic structural kernel, the following implication chain holds:
\[
\mathrm{GroundLinked}(x)\wedge \mathrm{FrameShifted}(x)
\Longrightarrow \mathrm{HCForced}(x),
\]
\[
\mathrm{Triolic}(x)\wedge \mathrm{Anchored}(x)\wedge \mathrm{HCForced}(x)
\Longrightarrow \mathrm{ClosureStable}(x),
\]
\[
\mathrm{ClosureStable}(x)\wedge \mathrm{HigherConfined}(x)
\Longrightarrow \mathrm{SeptinionReadable}(x).
\]
\end{proposition}

\begin{proof}
\textbf{First implication:
$\mathrm{GroundLinked}(x)\wedge\mathrm{FrameShifted}(x)
\Rightarrow\mathrm{HCForced}(x)$.}

$\mathrm{GroundLinked}(x)$ means $x\in S_n$ for some stratum $n$; it is
reachable by the cascade $\hat{K}$ from Definition~\ref{def:semantic-structural-kernel}~(S1).
$\mathrm{FrameShifted}(x)$ means by~(S8) that $x$ is a shifted or
regime-dependent readout of a common structural body.
By Lemma~\ref{lem:one-body-lemma}, the shifted presentations all refer
to one and the same carrier. By Lemma~\ref{lem:re-identification-lemma},
admissibility then requires a re-identification mechanism returning
competing presentations to a common output channel. The only operator
in the semantic kernel with this role is the compensator $C(x)=\tfrac{1}{2}(x+J(x))$
from~(S5), which is the unique $J$-symmetric projection.
Hence $\mathrm{HCForced}(x)$: the compensator is structurally necessary.

\smallskip
\textbf{Second implication:
$\mathrm{Triolic}(x)\wedge\mathrm{Anchored}(x)\wedge\mathrm{HCForced}(x)
\Rightarrow\mathrm{ClosureStable}(x)$.}

This is Lemma~\ref{lem:closure-forcing-lemma}, now with an explicit
four-step proof citing
Lemmas~\ref{lem:triolic-minimality-lemma},~\ref{lem:orthogonal-anchor-lemma},
\ref{lem:hc-admissibility-lemma},~\ref{lem:hc-idempotence-lemma},
and Theorem~\ref{thm:fractal}.

\smallskip
\textbf{Third implication:
$\mathrm{ClosureStable}(x)\wedge\mathrm{HigherConfined}(x)
\Rightarrow\mathrm{SeptinionReadable}(x)$.}

$\mathrm{ClosureStable}(x)$ means $\hat{K}(\mathcal{C}x)=\mathcal{C}x$
(Lemma~\ref{lem:closure-stable-iff-im-C}):
$x\in\mathrm{im}(C)$ is already in the ground sector.
$\mathrm{HigherConfined}(x)$ means $x$ is contained in a higher coupled
algebraic container — specifically the septinion/seption embedding
(Lemma~\ref{lem:higher-confinement-lemma}).
Within the higher container, the closure-stable state is readable as a
seption element precisely because its compensator-stabilized image
$\mathcal{C}x$ already lies in the admissible ground sector $S_0$, and
the higher container adds no new ontology
(Lemma~\ref{lem:non-escape-lemma}: no admissible transition exports $x$
into an independent external sector).
Hence the higher-confinement reading is a restatement of the already
established closure at higher algebraic level:
$\mathrm{SeptinionReadable}(x)$.
\end{proof}

\begin{proposition}[Tensorial intertwining conditions for the Millennium matrix]
\label{prop:tensorial-intertwining-m}
Every admissible tensor realization in the sense of
Definition~\ref{def:tensor-structural-kernel} satisfies:
\[
\mathcal M\mathcal J=\mathcal J\mathcal M,
\qquad
\mathcal M\mathcal C=\mathcal C\mathcal M,
\]
and, whenever layer transport is required to preserve the cascade reading,
\[
\mathcal M\hat K=\hat K\mathcal M.
\]
\end{proposition}

\begin{proof}
These are precisely condition~(T4) of
Definition~\ref{def:tensor-structural-kernel}, written in operator
notation. A tensor realization is admissible by definition only if
$\mathcal M$ satisfies the intertwining identities~(T4); hence the
relations hold for every admissible realization. No additional argument
is required beyond the definition.

The structural motivation for imposing~(T4) is as follows. The first
relation $\mathcal M\mathcal J=\mathcal J\mathcal M$ ensures that
transport by $\mathcal M$ does not destroy the involutive dual/frame
symmetry encoded by $\mathcal J$~(T1). The second relation
$\mathcal M\mathcal C=\mathcal C\mathcal M$ ensures that the
compensator-selected admissible sector $\mathrm{im}(\mathcal C)$ is
$\mathcal M$-invariant, consistent with $\mathcal C$ being idempotent
(T2). The third relation, when imposed, says that layer conversion
commutes with the cascade readout~(T3). These three motivations
explain why~(T4) is the correct admissibility condition; they do not
constitute an independent proof of~(T4) from (T1)--(T3).
\end{proof}

\begin{remark}[Resolution of Roadmap item A2]
\label{rem:a2-resolved}
Roadmap item~A2 identified that the intertwining relation
$\mathcal M\mathcal C=\mathcal C\mathcal M$ was supported only by a
necessity argument rather than a derivation. The resolution adopted
in v2.27.6 is to promote~(T4) to an explicit definitional condition
on admissible tensor realizations
(Definition~\ref{def:tensor-structural-kernel}, condition~(T4)).
Proposition~\ref{prop:tensorial-intertwining-m} then follows by
definition. The conditional marker on
Theorem~\ref{thm:hc-spectral-projector}
(Remark~\ref{rem:hc-spectral-projector-conditional}) is therefore
resolved: the theorem is no longer conditional on an unproven
hypothesis, but on an explicit definitional requirement.
\end{remark}

\begin{remark}[Semantic--tensor correspondence]
\label{rem:semantic-tensor-correspondence}
The semantic predicates and the tensor notation are intended to be read in
parallel. In particular,
\[
\mathrm{HCForced}(x)
\quad\leftrightarrow\quad
x\mapsto \mathcal Cx,
\]
\[
\mathrm{ClosureStable}(x)
\quad\leftrightarrow\quad
\hat K(\mathcal Cx)=\mathcal Cx,
\]
and
\[
\mathrm{Admissible}(x)
\quad\Rightarrow\quad
\mathrm{Admissible}(\mathcal Mx)
\]
whenever $\mathcal M$ satisfies the intertwining conditions of
Proposition~\ref{prop:tensorial-intertwining-m}. This makes it possible to
write later structure conditions in the same tensorial style as the
Riemann/Sphaera core while keeping their semantic meaning explicit.
\end{remark}


\subsection*{Phase I: First Unfolding Lemmas for the Generative Layer}
\addcontentsline{toc}{subsection}{Phase I: First Unfolding Lemmas for the Generative Layer}

We now begin the actual mathematical unfolding of the normalized structural
core by isolating the first chain of generative lemmas. The aim is not yet to
expand the full field/probe realization, but to make explicit the initial
forcing route from a common structural body to triolic organization,
orthogonal anchoring, frame-shifted readout, re-identification, and finally the
necessity of the Heisenberg compensator.

\begin{lemma}[One-body lemma]
\label{lem:one-body-lemma}
The semantic and tensorial kernels are to be read as describing different
admissible presentations of one and the same structural body, not independent
ontological objects. In particular, the cascade hierarchy, dual/frame symmetry,
and compensator projection all act on a common carrier rather than on unrelated
state spaces.
\end{lemma}

\begin{proof}
In Definition~\ref{def:semantic-structural-kernel} the kernel data are given as
one coordinated family
\[
(\mathbb S,\{S_n\}_{n\ge 0},S_0,\hat K,J,C,\mathcal M),
\]
not as several unrelated structures. The grouped architecture of
Definition~\ref{def:grouped-structural-architecture} likewise states that the
generative, projection, and field/probe layers are different organizational
regimes of the same underlying body. The tensor realization of
Definition~\ref{def:tensor-structural-kernel} then presents the same body in
operator form. Hence the structural reading is one-body throughout.
\end{proof}

\begin{lemma}[Triolic minimality lemma]
\label{lem:triolic-minimality-lemma}
A merely binary decomposition is structurally insufficient for the normalized
core. The minimal nondegenerate organization compatible with admissible
selection, compensator action, and closure readout is triolic.
\end{lemma}

\begin{proof}
By Definition~\ref{def:semantic-structural-kernel}, a triolic state carries a
three-sector organization $(K_1,K_2,K_3)$, while the orthogonal anchor singles
out the third sector as stabilizing direction relative to the first two. If one
retained only a binary split, then duality between the first two sectors could
still be stated, but there would be no distinguished third channel in which the
compensator-selected output and closure/light readout could be structurally
stabilized. The normalized core therefore requires three sectors as the minimal
nondegenerate organization.
\end{proof}

\begin{lemma}[Orthogonal anchor lemma]
\label{lem:orthogonal-anchor-lemma}
Within a triolic organization, the third sector is not an optional decoration but
the structural anchor that prevents the first two sectors from remaining a
merely parallel or relabelable pair. In particular,
\[
K_3\perp K_1,
\qquad
K_3\perp K_2.
\]
\end{lemma}

\begin{proof}
The orthogonal anchor condition is built into
Definition~\ref{def:semantic-structural-kernel}. Its role is clarified by the
forcing chain of Proposition~\ref{prop:semantic-forcing-chain}: closure
stability does not follow from duality alone, but from triolic structure plus a
protected anchor together with HC forcing. The third sector therefore fixes the
otherwise ambiguous two-sector pairing and provides the admissible direction in
which collapse/light selection can be stabilized.
\end{proof}

\begin{lemma}[Frame-shift lemma]
\label{lem:frame-shift-lemma}
Admissible effective states are read out in a phase- or frame-shifted
presentation. Accordingly, an effective state is not introduced as an
independent object but as a shifted view of the common structural body.
\end{lemma}

\begin{proof}
This is exactly the content of the frame-shift principle in
Definition~\ref{def:semantic-structural-kernel}, where
$\mathrm{FrameShifted}(x)$ means that $x$ is a shifted or regime-dependent
readout of a common body. The grouped architecture makes the same point at the
layer level: projection/translation and field/probe realizations are
organizational regimes of one underlying structure and therefore inherit rather
than replace the generative body.
\end{proof}

\begin{lemma}[Re-identification lemma]
\label{lem:re-identification-lemma}
If a state is simultaneously one-body and frame-shifted, then admissibility
requires a re-identification mechanism that returns competing presentations to a
common structural channel.
\end{lemma}

\begin{proof}
By Lemma~\ref{lem:one-body-lemma}, the relevant descriptions concern one and the
same body. By Lemma~\ref{lem:frame-shift-lemma}, these descriptions appear in
shifted or regime-dependent presentations. Without re-identification, the
shifted presentations would remain only formally related and no unique
admissible output sector could be selected. But the semantic kernel is built to
support admissible cascade readout and closure selection; hence a mechanism that
returns the shifted presentations to a common channel is structurally required.
\end{proof}

\begin{theorem}[HC-forcing theorem]
\label{thm:hc-forcing-theorem}
Within the normalized structural core, the required re-identification mechanism
is the Heisenberg compensator. Equivalently,
\[
\mathrm{GroundLinked}(x)\wedge \mathrm{FrameShifted}(x)
\Longrightarrow \mathrm{HCForced}(x).
\]
\end{theorem}

\begin{proof}
Proposition~\ref{prop:semantic-forcing-chain} already states the forcing
implication. The present point is to identify why the implication is necessary.
\[
\begin{aligned}
&\text{ground condition}
\Longrightarrow
\text{one triolic projection}
\Longrightarrow
\text{triolic minimality}
\Longrightarrow
\text{orthogonal anchor}
\\[0.2em]
&\Longrightarrow
\text{frame-shifted readout}
\Longrightarrow
\text{re-identification}
\Longrightarrow
\mathrm{HC}.
\end{aligned}
\]
\end{proof}

\begin{corollary}[First generative forcing chain]
\label{cor:first-generative-forcing-chain}
The first unfolded forcing route of the normalized core is
\[
\begin{aligned}
\text{one body}
&\Longrightarrow \text{triolic minimality}\\
&\Longrightarrow \text{orthogonal anchor}\\
&\Longrightarrow \text{frame-shifted readout}\\
&\Longrightarrow \text{re-identification}\\
&\Longrightarrow \mathrm{HC}.
\end{aligned}
\]
This chain is the entry point for the later closure and higher-confinement
unfolding.
\end{corollary}

\begin{proof}
Lemmas~\ref{lem:one-body-lemma}--\ref{lem:re-identification-lemma} establish the
successive structural requirements, and
Theorem~\ref{thm:hc-forcing-theorem} identifies their first nontrivial forced
operator consequence. The later closure and higher-confinement chain therefore
starts from this generative route.
\end{proof}


\subsection*{Phase I-B: HC and Closure Unfolding}
\addcontentsline{toc}{subsection}{Phase I-B: HC and Closure Unfolding}

We next unfold the first closure layer that begins once HC has been forced.
The aim of this block is to make precise that HC is not merely present, but
acts as an admissibility-preserving, anchor-compatible, closure-selecting
mechanism, and that exact closure can be distinguished from residual or only
partial closure.

\begin{lemma}[HC idempotence lemma]
\label{lem:hc-idempotence-lemma}
Once a state has been compensator-normalized, repeated compensator action does
not move it to a different admissible sector. Equivalently,
\[
C(C(x)) = C(x).
\]
\end{lemma}

\begin{proof}
By Definition~\ref{def:tensor-structural-kernel}, the tensorial compensator
satisfies $\mathcal C^2=\mathcal C$, and by
Definition~\ref{def:semantic-structural-kernel} the semantic compensator is the
corresponding map
\[
C(x)=\tfrac12\bigl(x+J(x)\bigr).
\]
Thus compensator action is projective rather than drifting: once a state lies
in the compensator-selected channel, further application leaves it there.
\end{proof}

\begin{lemma}[HC admissibility lemma]
\label{lem:hc-admissibility-lemma}
If $x$ is structurally admissible, then its compensator-normalized image $C(x)$
remains admissible.
\end{lemma}

\begin{proof}
The semantic/tensor kernel was introduced precisely so that the compensator is
not an external deformation but an internal normalization operator. Since
Theorem~\ref{thm:hc-forcing-theorem} identifies HC as the required
re-identification mechanism, using $C$ restores competing shifted presentations
to a common admissible channel instead of producing a foreign state. In the
tensor realization this is the admissible output selected by $\mathcal C$; on
the semantic side the same role is expressed by the preservation of the common
structural body under compensator projection. Hence admissibility is preserved.
\end{proof}

% ------------------------------------------------------------------
% Spectral-projector block (v2.27.4 integration)
% ------------------------------------------------------------------

\begin{definition}[Spectral decomposition of the state space]
\label{def:spectral-decomposition-M}
Let $\mathcal{M}$ denote the Millennium matrix acting on the admissible state
space $\mathcal{S}$. A \emph{spectral decomposition of $\mathcal{S}$ under
$\mathcal{M}$} is an orthogonal direct sum
\[
  \mathcal{S} \;=\; \bigoplus_{\lambda} \mathcal{S}_\lambda,
\]
where each $\mathcal{S}_\lambda \subset \mathcal{S}$ is a closed
$\mathcal{M}$-invariant subspace satisfying
$\mathcal{M}|_{\mathcal{S}_\lambda} = \lambda\,\mathrm{id}$.
A state $s \in \mathcal{S}$ is called \emph{spectrally stable} if
$s \in \mathcal{S}_\lambda$ for some $\lambda$, i.e.\ $\mathcal{M}(s) =
\lambda\,s$.
\end{definition}

\begin{remark}
Definition~\ref{def:spectral-decomposition-M} extends the existing
decomposition $\Hplus = \Hplusplus \oplus \Hplusminus$
(cf.\ Theorem~\ref{thm:unity}) to the full admissible readout regime.
The two primary spectral subspaces under any $\mathcal{M}$ satisfying
Proposition~\ref{prop:tensorial-intertwining-m} are precisely $\Hplusplus$
and $\Hplusminus$, since both are $\mathcal{M}$-invariant by the intertwining
condition $\mathcal{M}\mathcal{J} = \mathcal{J}\mathcal{M}$.
\end{remark}

\begin{lemma}[Diagonal self-traversal forces scalar action on the compensated sector]
\label{lem:diagonal-traversal-scalar-action}
Assume the Sphaera self-traversal condition of Theorem~\ref{thm:fractal}:
the cascade operator $\hat{K}$ transports the null axis $L_n$ of each
stratum $S_n$ to the null axis $L_{n-1}$ of the preceding stratum,
carrying the same structural content without generating new independent
directions. Then the Millennium matrix $\mathcal{M}$, when restricted to
the image of the compensator $\mathrm{im}(\mathcal{C})$, acts as a scalar
multiple of the identity:
\[
  \mathcal{M}\big|_{\mathrm{im}(\mathcal{C})} \;=\; \lambda_0\,\mathrm{id},
\]
where $\lambda_0$ is the unique admissible eigenvalue of $\mathcal{M}$ on
the closure-stable sector.
\end{lemma}

\begin{proof}
By Theorem~\ref{thm:fractal}, the self-traversal of the Sphaera is
\emph{diagonal}: the cascade does not generate a new structural object at
each stratum, but transports one and the same null-axis content through
admissible realization levels. In operator language this means that the
cascade direction at each stratum is not a Jordan-chain descendant of the
previous one---no new vector is generated from an existing vector by a
nilpotent shift.

By the intertwining condition $\mathcal{M}\mathcal{C} = \mathcal{C}\mathcal{M}$
(Proposition~\ref{prop:tensorial-intertwining-m}), $\mathrm{im}(\mathcal{C})$
is $\mathcal{M}$-invariant. Suppose for contradiction that
$\mathcal{M}|_{\mathrm{im}(\mathcal{C})}$ had a non-trivial Jordan block:
then there would exist $s,w \in \mathrm{im}(\mathcal{C})$ with $w\neq 0$ and
\[
  \mathcal{M}(s) = \lambda_0\,s + w.
\]
This would mean $\mathcal{M}$ generates $w$ from $s$---a new independent
direction within the compensated sector---contradicting the diagonal
self-traversal property. Hence $\mathcal{M}|_{\mathrm{im}(\mathcal{C})}$
is diagonalizable. Since $\mathrm{im}(\mathcal{C})$ is the closure-stable
sector (Lemma~\ref{lem:closure-forcing-lemma}), which is irreducible under
$\mathcal{M}$, a diagonalizable action on an irreducible invariant subspace
is scalar. Denoting this scalar $\lambda_0$ completes the proof.
\end{proof}

\begin{theorem}[HC as spectral projector of the Millennium matrix]
\label{thm:hc-spectral-projector}
Assume the admissible state space $\mathcal{S}$ carries a spectral
decomposition under $\mathcal{M}$ as in
Definition~\ref{def:spectral-decomposition-M}, that $\mathcal{M}$ satisfies
the intertwining conditions of Proposition~\ref{prop:tensorial-intertwining-m},
and that the Sphaera self-traversal is diagonal in the sense of
Lemma~\ref{lem:diagonal-traversal-scalar-action}. Then the Heisenberg
compensator $C$ coincides with the spectral projector of $\mathcal{M}$
onto the closure-stable eigenspace $\mathcal{S}_0 :=
\ker(\mathcal{M} - \lambda_0\,\mathrm{id})$:
\[
  C \;=\; \Pi^{(\mathcal{M})}_{\mathrm{spec}}.
\]
\end{theorem}

\begin{remark}[Status of Theorem~\ref{thm:hc-spectral-projector}]
\label{rem:hc-spectral-projector-conditional}
As of v2.27.6, both open premises of this theorem are resolved:
the bijection $\mathrm{ClosureStable}(x)\Leftrightarrow x\in\mathrm{im}(C)$
is established by Lemma~\ref{lem:closure-stable-iff-im-C} (Roadmap A3),
and the intertwining relation $\mathcal{M}\mathcal{C}=\mathcal{C}\mathcal{M}$
is now a definitional condition (T4) of
Definition~\ref{def:tensor-structural-kernel} (Roadmap A2,
Remark~\ref{rem:a2-resolved}).
The theorem therefore holds unconditionally within every admissible
tensor realization.
\end{remark}

\begin{proof}
We verify the three conditions characterizing a spectral projector.

\smallskip
\noindent\textbf{Step 1: Idempotence.}
By Lemma~\ref{lem:hc-idempotence-lemma}, $C^2=C$. By
Definition~\ref{def:tensor-structural-kernel},
$\mathcal{C}^\mu{}_\alpha\mathcal{C}^\alpha{}_\nu = \mathcal{C}^\mu{}_\nu$
in the tensor realization. Hence $C$ is a projector.

\smallskip
\noindent\textbf{Step 2: $\mathcal{M}$-compatibility.}
By Proposition~\ref{prop:tensorial-intertwining-m},
$\mathcal{M}\mathcal{C} = \mathcal{C}\mathcal{M}$. Hence $\mathrm{im}(C)$
is $\mathcal{M}$-invariant and $C$ projects onto a $\mathcal{M}$-stable
subspace.

\smallskip
\noindent\textbf{Step 3: Scalar eigenvalue on the image.}
By Lemma~\ref{lem:diagonal-traversal-scalar-action}, the diagonal
self-traversal forces $\mathcal{M}|_{\mathrm{im}(C)} = \lambda_0\,\mathrm{id}$.
Hence $\mathrm{im}(C) \subseteq \mathcal{S}_0$.
By Lemma~\ref{lem:closure-forcing-lemma}, the closure-stable states are
exactly those in $\mathrm{im}(C)$, so $\mathrm{im}(C) = \mathcal{S}_0$.

\smallskip
\noindent\textbf{Conclusion.}
A projector with $\mathcal{M}$-invariant image equal to a single eigenspace
$\mathcal{S}_0$ is the spectral projector $\Pi^{(\mathcal{M})}_{\mathrm{spec}}$
onto that eigenspace.
\end{proof}

\begin{corollary}[Structural fixpoint: $\mathcal{M}$ and $C$ mutually determine each other]
\label{cor:M-HC-fixpoint}
Under the hypotheses of Theorem~\ref{thm:hc-spectral-projector},
\[
  \mathcal{M}\,C \;=\; C\,\mathcal{M} \;=\; \lambda_0\,C.
\]
\end{corollary}

\begin{proof}
For any $s\in\mathcal{S}$: since $C(s)\in\mathcal{S}_0$, we have
$\mathcal{M}(C(s))=\lambda_0\,C(s)$, giving $\mathcal{M}\,C=\lambda_0\,C$.
For $s\in\mathcal{S}_0$: $C\,\mathcal{M}(s)=C(\lambda_0 s)=\lambda_0 C(s)$.
For $s\in\ker(C)$: $C\,\mathcal{M}(s)=0$ (since $\ker(C)$ is
$\mathcal{M}$-invariant and maps to a sector annihilated by $C$).
Hence $C\,\mathcal{M}=\lambda_0\,C$.
\end{proof}

\begin{remark}[Phase-shifted frame and spectral measurement]
\label{rem:frame-shift-spectral}
The identification $C = \Pi^{(\mathcal{M})}_{\mathrm{spec}}$ has a direct
consequence for the matter frame. Since matter observers reside at
$\zminus = +\bsym \approx 0.5493$ rather than the photon ground shell
$\zminus = 0$, every spectral readout of $\mathcal{M}$ is performed in the
shifted frame. The measured eigenvalue is therefore not $\lambda_0$ itself
but the frame-corrected value
\[
  \lambda_0^{(\mathrm{matter})}
  \;=\; \gL^{-k(\lambda_0)}\cdot\lambda_0^{(\mathrm{photon})},
\]
where $k(\lambda_0)$ is the frame-insertion count of the observable and
$\gL \approx 1.864$ is the matter--photon Lorentz factor
(Theorem~\ref{thm:frame-offset}). All physical constants derived from
spectral eigenvalues of $\mathcal{M}$ therefore carry systematic
$\gL^k$-corrections relative to their photon-frame values.
\end{remark}

\begin{corollary}[Standard Model realization of the spectral decomposition]
\label{cor:SM-spectral-realization}
The octonionic decomposition $\OO = \mathbb{C}\oplus\mathbb{C}^3$
(Theorem~\ref{thm:octonion-SU3}) is a concrete realization of
Definition~\ref{def:spectral-decomposition-M}: the singlet $\mathbb{C}$ is
the closure-stable eigenspace $\mathcal{S}_0 = \mathrm{im}(C)$, and the
triplet $\mathbb{C}^3$ carries the three color eigenspaces
$\mathcal{S}_{K_1},\mathcal{S}_{K_2},\mathcal{S}_{K_3}$.
Under this identification, Theorem~\ref{thm:hc-spectral-projector}
specializes as follows:
\begin{enumerate}[label=(\roman*)]
  \item The compensator $C$ projects onto the singlet $\mathbb{C}\subset\OO$,
  which is the photon kernel $K_3$ sector ($\zminus=0$).
  \item The triplet $\mathbb{C}^3$ is $\mathcal{M}$-invariant and carries
  the SU(3) color action; matter ($K_1$) and tachyon ($K_2$) sectors are
  spectral subspaces at non-zero frame offset.
  \item The matter-frame correction reads
  \[
    \lambda_0^{(\mathrm{matter})}
    = \gL^{-k(\lambda_0)}\cdot\lambda_0^{(\mathrm{photon})},
  \]
  consistent with Theorem~\ref{thm:SM}: $K_3$ is frame-invariant
  ($k=0$), while the triplet components carry $k\ge 1$ corrections
  and are therefore observed with $\gL^k$-suppression in the matter frame.
\end{enumerate}
\end{corollary}

\begin{proof}
By Theorem~\ref{thm:octonion-SU3}, the octonion automorphism group $G_2$
restricts to $\mathrm{SU}(3)\subset G_2$, yielding the direct sum
$\OO = \mathbb{C}\oplus\mathbb{C}^3$. The singlet $\mathbb{C}$ is the
unique $\mathrm{SU}(3)$-invariant subspace, hence $\mathcal{M}$-stable at
a scalar eigenvalue $\lambda_0$. By the remark in
Theorem~\ref{thm:SM}, $K_3 = e_1 e_2$ is simultaneously the photon kernel
and the projection of the Heisenberg compensator; hence
$\mathrm{im}(C) = \mathbb{C} = \mathcal{S}_0$. The three kernels
$(K_1,K_2,K_3)$ realize the three colors of $\mathbb{C}^3$, confirming
$\mathcal{S}_{K_i}$ as spectral subspaces. Statement~(iii) then follows
directly from Remark~\ref{rem:frame-shift-spectral} applied to each
triplet component.
\end{proof}

% ------------------------------------------------------------------
% Readout-Relative Lemmas and Structural Lifts (v2.27.4 integration)
% ------------------------------------------------------------------

\section{Readout-Relative Lemmas and Structural Lifts}
\label{sec:readout-relative-lemmas}

The present framework distinguishes sharply between the structural kernel
of the theory and the visible form in which that kernel is read in a
given regime or frame. This distinction is not optional: the same
closure-stable object may appear as a probe/Dirac-type law, as an
effective QFT-type law in a phase-shifted observer regime, or as an
Einstein-type law after large-scale geometric readout. Consequently,
one must distinguish between \emph{kernel-level truths} and
\emph{readout-relative truths}.

The concrete instantiation of this distinction is already present in the
preceding spectral-projector block: the eigenvalue $\lambda_0$ is a
kernel-level invariant, while the measured value
$\lambda_0^{(\mathrm{matter})} = \gL^{-k}\cdot\lambda_0^{(\mathrm{photon})}$
(Remark~\ref{rem:frame-shift-spectral}) is readout-relative, depending
on the frame-insertion count $k$ of the observable in the matter frame
at $\zminus = +\bsym$.

\subsection*{Kernel level and readout level}

\begin{definition}[Kernel-level statement]
\label{def:kernel-level-statement}
A statement is called \emph{kernel-level} if it is invariant under the
admissible structural equivalences of the framework, including:
\begin{enumerate}[label=(\roman*)]
  \item sector-preserving unitary equivalence;
  \item compensator projection;
  \item observer-frame phase shift;
  \item admissible branch restriction.
\end{enumerate}
Such statements express properties of the underlying closure-stable
structural object itself rather than of one distinguished reading of it.
\end{definition}

\begin{definition}[Readout-relative lemma]
\label{def:readout-relative-lemma}
A statement $L^{(f)}$ is called a \emph{readout-relative lemma} if it
is proven only in a distinguished regime, frame, or readout $f$, and its
visible form depends on that readout. In particular, a readout-relative
lemma need not remain true as a frame-blind absolute statement.
\end{definition}

\begin{remark}
This is not a defect of the method. It is the correct formal response to
the fact that different regime-readings expose different visible shadows
of one common kernel. In the present framework this applies in particular
to the probe, effective QFT, and Einstein readouts.
The Symmetric matter--tachyon frame theorem (Theorem~\ref{thm:symmetric-frame})
is an example: $\beta_{\mathrm{matter}} = +\bsym$ is readout-relative
(it is the matter-frame reading), while the symmetry
$\beta_{\mathrm{tachyon}} = -\bsym$ and the common kernel $\bsym$ are
kernel-level invariants.
\end{remark}

\subsection*{Why this distinction is necessary}

The need for readout-relative lemmas is already built into the
architecture. The master system is closure-stable and unique only up to
the allowed sector-preserving, unitary, and frame transformations, while
the effective regimes differ by compensator decomposition, large-scale
suppression of closure fluctuations, and geometric readout. Thus one and
the same kernel may generate visibly different but structurally
compatible statements. In the present release this distinction should be read together with the OP transport theorem and the visible/structural mismatch separation tool: visible disagreement is secondary unless native closure-class preservation itself fails.

\begin{proposition}[Legitimacy of restricted lemmas]
\label{prop:legitimacy-restricted-lemmas}
A lemma proven in a distinguished readout $f$ may be used in the general
argument provided that:
\begin{enumerate}[label=(\roman*)]
  \item its frame or regime of validity is stated explicitly;
  \item its invariant kernel content is separated from its visible
        readout form;
  \item the passage from the readout-relative statement to the general
        claim is justified by a structural transport principle.
\end{enumerate}
\end{proposition}

\begin{proof}
The framework already distinguishes the common structural equation from
its regime projections (Theorem~\ref{thm:unity}, unitary equivalence via
Stone--von Neumann). Therefore a proof in one readout need not be
discarded merely because its visible form changes in another readout.
What must be transported is not the literal visible formula, but the
kernel-level content preserved under the admissible transport operations
of Lemmas~\ref{lem:group-transport} and~\ref{lem:hc-admissibility-lemma}.
\end{proof}

\subsection*{Structural lifts}

\begin{definition}[Structural lift]
\label{def:structural-lift}
Let $L^{(f)}$ be a readout-relative lemma valid in readout $f$.
A \emph{structural lift} of $L^{(f)}$ is a kernel-level statement
$L_{\mathrm{struct}}$ obtained by extracting from $L^{(f)}$ only the
invariant content that survives admissible unitary, sector, and frame
transport.
\end{definition}

\begin{remark}
A structural lift is therefore weaker than a literal universalization of
$L^{(f)}$, but stronger than leaving the lemma imprisoned in one special
case. The corollary Cor.~\ref{cor:SM-spectral-realization} is a concrete
example: statement~(iii) there lifts the frame-corrected eigenvalue
$\lambda_0^{(\mathrm{matter})}$ to the kernel-level invariant $\lambda_0$
by identifying $k=0$ for the singlet and $k\ge 1$ for the triplet.
\end{remark}

\begin{proposition}[Structural lift principle]
\label{prop:structural-lift-principle}
Suppose $L^{(f)}$ is proven in readout $f$, and suppose its proof
depends only on:
\begin{enumerate}[label=(\roman*)]
  \item the closure-stable master structure;
  \item invariants preserved under admissible transport;
  \item readout-specific visibility factors that can be isolated.
\end{enumerate}
Then there exists a structural lift $L_{\mathrm{struct}}$ that is valid
at kernel level, even if the visible statement $L^{(f)}$ itself is not
frame-blindly universal.
\end{proposition}

\begin{proof}
\noindent\textbf{Invariant content is transport-stable.}
By hypotheses~(i) and~(ii), the kernel-invariant content of the proof
of $L^{(f)}$ is preserved by the admissible transport operations.
Concretely: Lemma~\ref{lem:group-transport} guarantees that grouped-layer
transport carries admissible structure without creating new ontology, and
Lemma~\ref{lem:hc-admissibility-lemma} guarantees that compensator action
leaves the admissible sector invariant. Hence the kernel-invariant content
$P_{\mathrm{inv}}(\text{proof of }L^{(f)})$ is valid in every admissible
readout, not only in $f$.

\smallskip
\noindent\textbf{Readout-specific factors are isolatable.}
By hypothesis~(iii), the readout-dependent terms form a separable complement.
Discarding them while retaining $P_{\mathrm{inv}}(\text{proof})$ yields
$L_{\mathrm{struct}}$, independent of $f$.

\smallskip
\noindent\textbf{$L_{\mathrm{struct}}$ is kernel-level.}
Since $L_{\mathrm{struct}}$ depends only on the closure-stable structure
and transport-preserved invariants, it satisfies all four conditions of
Definition~\ref{def:kernel-level-statement}: invariance under
sector-preserving unitary equivalence, compensator projection,
observer-frame phase shift, and admissible branch restriction.
\end{proof}

\subsection*{The three-step proof logic}

This principle simplifies the global architecture considerably.
Instead of demanding that every intermediate statement already be valid
as a frame-blind absolute theorem, one may work in three steps:
\[
  \text{readout-relative lemma}
  \;\Longrightarrow\;
  \text{structural lift}
  \;\Longrightarrow\;
  \text{general kernel-level theorem}.
\]
This is the natural proof logic for a theory whose visible regimes are
different projections of one and the same structural object.

\subsection*{Shadow versus core: consequence for the compensator}

\begin{remark}[HC shadow versus HC kernel]
\label{rem:hc-shadow-vs-core}
This distinction is particularly important for the compensator.
In the minimal formulation, the visible compensator functional
$\mathcal{H}_C^{(f)}$ behaves as a closure-mismatch penalty in a given
readout and may vanish in an exact sector ($\mathcal{H}_C = 0$ in the
photon frame; cf.\ the closed master system). But under the
frame-dependent readout hypothesis developed above, such vanishing
concerns only the \emph{visible defect shadow} in a given regime, not
the deeper compensator-generating kernel
$C = \Pi^{(\mathcal{M})}_{\mathrm{spec}}$
(Theorem~\ref{thm:hc-spectral-projector}).

Hence a statement about $\mathcal{H}_C^{(f)}$ must not automatically be
read as a statement about the entire HC architecture. In particular, the
vanishing of the visible compensator field in the photon-frame readout
is consistent with the HC kernel remaining active as the
spectral-projector mechanism that selects the closure-stable sector ---
it is the shadow, not the core, that vanishes.

This is directly relevant to OP~5 ($\Lambda_{\mathrm{obs}}$): the
cosmological constant appears as a \emph{readout-relative} quantity ---
the frame offset $\bsym^2\langle f,f\rangle$ visible in the matter frame
--- while the photon-frame value $\Lambda_{\mathrm{fundamental}} = 0$
is the kernel-level statement. The large observed value is not a
fine-tuning problem but a failure to separate readout-relative visibility
from structural invariance.
\end{remark}

% ------------------------------------------------------------------
% Bidirectional Probe Lensing (v2.27.4 integration)
% ------------------------------------------------------------------

\section{Bidirectional Probe Lensing Hypothesis}
\label{sec:bidirectional-probe-lensing}

The preceding sections showed that the Dirac-, QFT-, and Einstein-type
regimes are not separate theories, but structural projections of one
common closure-stable master equation
(Proposition~\ref{prop:structural-lift-principle}). They also showed
that the Einstein reconstruction presently depends on a bilinear coframe
extraction whose nondegeneracy is still model-dependent. This suggests a
different route: rather than attempting to reconstruct geometry from a
single readout only, one may compare the traces of two probes approaching
the same structural core from different regime readings.

\subsection*{Basic setup}

Consider two probe modes of the triadic system,
$\Psi_{\mathrm{ART}}$ and $\Psi_{\mathrm{QFT}}$,
together with the orthogonal triadic component $\Psi_\perp$.
We regard $\Psi_{\mathrm{ART}}$ as a probe emphasizing the
large-scale/geometric regime and $\Psi_{\mathrm{QFT}}$ as a probe
emphasizing the phase-shifted observer/QFT regime at
$\zminus = +\bsym$ (Theorem~\ref{thm:symmetric-frame}).
The combined field is treated as a specialization of the triadic
two-probe system
\[
  \Psi_{\triangle}
  =
  \begin{pmatrix}
    \Psi_{\mathrm{ART}}\\
    \Psi_{\mathrm{QFT}}\\
    \Psi_\perp
  \end{pmatrix},
  \qquad
  \bigl(i\Gamma^\mu\mathcal{D}_\mu - \Omega_\triangle\bigr)
  \Psi_\triangle = 0.
\]
The third component is not decorative: it is the intrinsic coupling
channel through which the two regime-readings communicate.

\subsection*{Structural hypothesis}

\begin{proposition}[Bidirectional probe lensing hypothesis]
\label{hyp:bidirectional-probe-lensing}
\emph{(Hypothesis-grade; not yet theorem-level.)}
Let two probes approach the same closure-stable structural core from
different regime-readings, one on the ART side and one on the QFT side.
Then the compensator/closure core acts as an effective lensing kernel
for their traces. More precisely:
\begin{enumerate}[label=(\roman*)]
  \item the two probes need not follow identical local trajectories;
  \item their relative deflection, phase lag, and focusing structure are
        controlled by the same compensator/closure sector
        $C = \Pi^{(\mathcal{M})}_{\mathrm{spec}}$
        (Theorem~\ref{thm:hc-spectral-projector});
  \item the orthogonal triadic channel $\Psi_\perp$ mediates the
        nontrivial coupling between the two traces;
  \item the resulting trace geometry carries information about the
        effective metric readout before a full bilinear coframe
        reconstruction is imposed.
\end{enumerate}
This is a readout-relative statement
(Definition~\ref{def:readout-relative-lemma}): the lensing geometry is
visible in the bidirectional probe readout, while the kernel-level
content is the common closure-stable core.
\end{proposition}

\begin{theorem}[Theorem-level shadow of the bidirectional probe lensing hypothesis]
\label{thm:theorem-level-shadow-bidirectional-probe-lensing}
At the active theorem level, the present framework already supports the following weaker probe-lensing shadow:
\begin{enumerate}[leftmargin=2em,label=(\roman*)]
  \item two probe readouts approaching the same closure-stable structural core may be treated as distinct regime-readings of one common carrier rather than as probes of two disconnected foundations;
  \item the orthogonal triadic component is already a genuine coupling channel between the two regime-readings, not a decorative extra field;
  \item what is \emph{not} yet theorem-level is a full geometric law that identifies the resulting relative deflection, focusing, or phase-lag pattern with one explicit effective lensing tensor in complete generality.
\end{enumerate}
Therefore the proven theorem-level shadow of Hypothesis~\ref{hyp:bidirectional-probe-lensing} is:
\emph{the two probes already share one closure-stable carrier and one intrinsic coupling channel, even though a full lensing law for their relative geometry is not yet proved.}
\end{theorem}

\begin{proof}
The active framework already treats the ART- and QFT-side equations as distinguished regime/readout projections of one common closure-stable structural kernel rather than as disconnected sector theories. In the bidirectional probe setup, the combined triadic field
\[
  \Psi_{\triangle}
  =
  \begin{pmatrix}
    \Psi_{\mathrm{ART}}\\
    \Psi_{\mathrm{QFT}}\\
    \Psi_\perp
  \end{pmatrix}
\]
is explicitly organized so that the orthogonal component $\Psi_\perp$ is the intrinsic channel through which the two regime-readings communicate. Hence the theorem-level content already available is the existence of one common carrier and one nontrivial coupling channel between the two probes. What remains unproved is the stronger claim that their relative deflection/focusing/phase-lag pattern is already governed by one explicit universal lensing law at theorem level.
\end{proof}

\begin{corollary}[Operational reading of the probe-lensing hypothesis in the present v\docversion\ release line]
\label{cor:operational-reading-bidirectional-probe-lensing-v320}
Until a full effective lensing law is derived, later arguments may already use Theorem~\ref{thm:theorem-level-shadow-bidirectional-probe-lensing} as the active theorem-level substitute for the weakest part of Hypothesis~\ref{hyp:bidirectional-probe-lensing}: one may work with two probe readouts of one common closure-stable carrier coupled through the orthogonal triadic channel, while keeping the stronger geometric lensing law itself deferred. In the v3.20 crystal language, this means that two probe-facing domain realizations may already be treated as role-faces of one common carrier corridor, without claiming a complete universal focusing tensor.
\end{corollary}

\begin{proof}
This is the operational reformulation of Theorem~\ref{thm:theorem-level-shadow-bidirectional-probe-lensing}. The theorem already isolates the precise point at which common-carrier/two-readout coupling is available theorematically, while the stronger lensing law remains hypothesis-grade.
\end{proof}

\begin{remark}[Status of the compressed interface-to-optical bridge layer in the present release line]
\label{rem:status-of-the-first-bundled-theorem-level-shadow-package-v320}
\label{rem:status-of-the-first-v320-crystal-shadow-interface}
\label{rem:status-of-the-terminal-v320-working-interface-surface}
\label{rem:status-of-the-first-interface-level-defect-face-vanishing-step}
\label{rem:status-of-the-first-interface-level-fluid-readout-demotion-step}
\label{rem:status-of-the-first-interface-level-common-carrier-projection-step}
\label{rem:status-of-the-first-interface-level-probe-coupling-step}
\label{rem:status-of-the-first-bundled-interface-application-package}
\label{rem:status-of-the-first-bundled-interface-compatibility-step-in-v321}
\label{rem:status-of-the-first-op3-op5-bridge-test-on-the-terminal-interface}
\label{rem:status-of-the-first-no-second-carrier-step-in-v322}
\label{rem:status-of-the-first-partial-op3-op5-closure-theorem-in-v322}
\label{rem:status-of-the-first-restricted-lift-back-step-in-v322}
\label{rem:status-of-the-first-prism-readout-reduction-step-in-v323}
\label{rem:status-of-the-first-residual-spectral-diagnostic-step-in-v323}
\label{rem:status-of-the-first-parametric-spectral-fit-layer-in-v323}
\label{rem:status-of-the-first-spectral-calibration-step-in-v323}
\label{rem:status-of-the-first-overdetermined-spectral-consistency-step-in-v323}
\label{rem:status-of-the-first-effective-index-reconstruction-step-in-v323}
\label{rem:status-of-the-first-carrier-profile-correction-step-in-v323}
\label{rem:status-of-the-first-reduced-spectral-stability-step-in-v323}
\label{rem:status-of-the-first-bundled-optical-approximation-package-in-v323}
The present line no longer expands the long v3.20--v3.23 ladder step by step inside the active core. Instead, the theorem-level shadow package, the crystal--shadow/interface layer, the terminal interface application package, the first OP~3/OP~5 bridge/closure consequences, and the first optical approximation package are retained through one compressed interface-to-optical bridge layer. The detailed development remains recoverable from earlier candidates, while the active line keeps only the citation-equivalent surface needed for current proof use.
\end{remark}

\begin{definition}[Compressed interface-to-optical bridge token]
\label{def:v320-crystal-shadow-interface-object}
\label{def:reduced-prism-readout-sector}
\label{def:prism-readout-residual-profile}
\label{def:first-effective-carrier-dispersion-fit-ansatz}
\label{def:first-reconstructed-effective-refractive-profile}
\label{def:first-carrier-profile-correction-decomposition}
Let
\[
\mathfrak B^{\mathrm{int\text{-}opt}}
:=
\bigl(
\Sigma^{\mathrm{shadow}}_{\mathrm{v20}},
\mathfrak I_{\mathrm{term}},
\mathfrak R_{\mathrm{prism}},
\rho_{\mathrm{res}},
\mathcal F_{\mathrm{disp}},
 n^{\mathrm{rec}}_{\mathrm{eff}},
 \Delta n_{\mathrm{corr}}
\bigr)
\]
denote the compressed interface-to-optical bridge token, where the entries stand respectively for the bundled theorem-level shadow package, the terminal crystal--shadow interface object, the reduced prism-readout sector, the residual spectral profile, the first effective dispersion-fit ansatz, the reconstructed effective index profile, and its first carrier-profile correction decomposition. The active line uses this token as the compact carrier of all v3.20--v3.23 interface-to-optical data that are still needed in the present v\docversion\ release line.
\end{definition}

\begin{lemma}[Compressed spectral entry lemma on the interface-to-optical bridge]
\label{lem:first-prism-readout-lemma-on-the-retained-carrier-surface}
\label{lem:first-residual-spectral-diagnostic-lemma}
From the token $\mathfrak B^{\mathrm{int\text{-}opt}}$ one may recover, up to citation-equivalent proof use, the reduced prism sector together with the first residual spectral diagnostic entry data. In particular, any later argument that needs only the passage from terminal interface readout to residual spectral comparison may cite the present lemma instead of reopening the separate prism-readout and residual-profile lemmas.
\end{lemma}

\begin{proof}
The token $\mathfrak B^{\mathrm{int\text{-}opt}}$ retains exactly the interface-level theorematic shadow content, the terminal crystal/interface content, and the first optical approximation data that were previously introduced as separate stages. Since the active line uses these stages only through their retained terminal/citation-equivalent content, the reduced prism entry and the first residual spectral comparison layer may be read back from the single token without reopening the former intermediate exposition.
\end{proof}

\begin{proposition}[Compressed interface-to-optical bridge package]
\label{prop:first-bundled-theorem-level-shadow-package}
\label{prop:first-v320-crystal-shadow-compatibility-principle}
\label{prop:terminal-citation-rule-for-the-present-v320-working-interface}
\label{prop:first-interface-level-defect-face-vanishing-principle}
\label{prop:first-interface-level-fluid-readout-demotion-principle}
\label{prop:first-interface-level-common-carrier-projection-principle}
\label{prop:first-interface-level-probe-coupling-principle}
\label{prop:first-bundled-interface-application-package}
\label{prop:first-op3-op5-bridge-test-on-the-terminal-interface}
\label{prop:first-parametric-spectral-fit-principle}
\label{prop:first-three-parameter-spectral-calibration-principle}
\label{prop:first-overdetermined-spectral-consistency-principle}
\label{prop:first-effective-index-reconstruction-principle-from-prism-data}
\label{prop:first-carrier-profile-correction-principle}
\label{prop:first-reduced-spectral-stability-principle}
\label{prop:first-bundled-optical-approximation-package}
Assume a later argument depends only on the theorem-level shadow/common-carrier content, the terminal interface/crystal application layer, the weak OP~3/OP~5 bridge consequences already retained on that terminal layer, and the first optical approximation package given by reduced prism readout, residual spectral comparison, parametric fit, calibration, consistency, effective-index reconstruction, carrier-profile correction, and reduced spectral stability. Assume moreover that the argument does not require any still-deferred stronger claim such as an explicit HC meta-field equation, a deeper compensator generator, a second primary carrier, a full universal probe-lensing law, or a stronger optical inversion theorem beyond the first retained approximation package. Then the argument may work entirely with the single bridge token $\mathfrak B^{\mathrm{int\text{-}opt}}$ instead of reopening the former v3.20--v3.23 ladder step by step.
\end{proposition}

\begin{proof}
All formerly separate ingredients in the stated range were already used in the active line only through retained theorem-level or citation-level content: bundled theorematic shadows, terminal interface application, weak OP~3/OP~5 bridge consequences, and the first optical approximation package. The token $\mathfrak B^{\mathrm{int\text{-}opt}}$ was defined precisely to retain this active content and forget only the now-redundant intermediate exposition. Therefore any later argument that uses no stronger claim than the ones listed can cite the bridge token and stay on the compressed interface-to-optical surface.
\end{proof}

\begin{theorem}[Compressed theorematic consequences on the interface-to-optical bridge]
\label{thm:first-bundled-interface-compatibility-theorem}
\label{thm:first-no-second-carrier-theorem-on-the-terminal-interface}
\label{thm:first-partial-op3-op5-closure-theorem-on-the-terminal-interface}
\label{thm:first-restricted-lift-back-theorem-from-the-terminal-interface}
On the admissible theorematic window represented by $\mathfrak B^{\mathrm{int\text{-}opt}}$, the active line already retains, in compressed form, the compatibility theorem of the terminal interface package, the no-second-carrier consequence, the first partial OP~3/OP~5 closure consequence, and the first restricted lift-back consequence. Hence these theorematic uses may be cited through the single compressed theorematic bridge whenever no finer reopening below the retained interface-to-optical surface is needed.
\end{theorem}

\begin{proof}
The separate theorems in the former expanded presentation all concerned theorematic consequences that had already stabilized on the same terminal interface/application layer and that were then carried forward only through their retained citation content. Since the token $\mathfrak B^{\mathrm{int\text{-}opt}}$ preserves precisely that retained layer, the corresponding theorematic consequences may likewise be cited in compressed form on the same admissible window.
\end{proof}

\begin{corollary}[Compressed use rule for the interface-to-optical bridge]
\label{cor:bundled-use-rule-v320-theorem-level-shadow-package}
\label{cor:use-rule-for-the-v320-crystal-shadow-interface}
\label{cor:operational-mini-checklist-for-the-terminal-v320-working-interface}
\label{cor:interface-level-use-rule-for-exact-sector-defect-face-vanishing}
\label{cor:interface-level-use-rule-for-coarse-grained-fluid-readout}
\label{cor:interface-level-use-rule-for-common-carrier-and-projection-architecture}
\label{cor:interface-level-use-rule-for-probe-coupled-readout-architecture}
\label{cor:bundled-use-rule-for-the-terminal-interface-application-package}
\label{cor:bundled-use-rule-for-the-first-interface-compatibility-theorem}
\label{cor:use-rule-for-the-first-op3-op5-bridge-test}
\label{cor:use-rule-for-the-first-no-second-carrier-theorem}
\label{cor:use-rule-for-the-first-partial-op3-op5-closure-theorem}
\label{cor:use-rule-for-the-first-restricted-lift-back-theorem}
\label{cor:first-spectral-approximation-entry-rule-from-prism-readout}
\label{cor:first-spectral-fit-entry-rule-from-prism-residuals}
\label{cor:first-carrier-fit-workflow-for-optical-data}
\label{cor:first-identifiability-rule-for-optical-fit-parameters}
\label{cor:first-validation-workflow-for-measured-optical-series}
\label{cor:first-inversion-workflow-from-prism-angles-to-carrier-profile}
\label{cor:first-refinement-workflow-for-the-effective-carrier-law}
\label{cor:first-stopping-rule-for-the-reduced-prism-readout-fit-loop}
\label{cor:bundled-use-rule-for-the-first-optical-approximation-package}
Whenever a later argument needs only the active content compressed into $\mathfrak B^{\mathrm{int\text{-}opt}}$, it is sufficient to cite the present bridge layer and proceed directly to the present optical-diagnostic control surface. A finer reopening is required only when the argument depends on information intentionally forgotten by the compression: stronger generator-level structure, a deeper carrier split, a stronger universal lensing law, or a more detailed optical inversion/construction below the retained first approximation package.
\end{corollary}

\begin{proof}
This is the operational reformulation of Proposition~\ref{prop:first-bundled-optical-approximation-package} and Theorem~\ref{thm:first-partial-op3-op5-closure-theorem-on-the-terminal-interface} in their compressed bridge form. The present layer isolates the exact point at which the active line may stay on the retained interface-to-optical surface and pass directly into the present optical-diagnostic control machinery.
\end{proof}

\begin{remark}[Status of the compressed optical-diagnostic control layer in the present release line]
\label{rem:status-of-the-compressed-optical-diagnostic-control-layer-in-v324}
The present v\docversion\ release line now carries the whole first optical-to-control pipeline in compressed form: reference comparison, shortlist formation, structural back-reading, reopening priority, bombe/sieve reduction, tensor/shell control, early-depth policy, local utility ranking, and the first local spectral/QAM readout are retained as one finite citation layer rather than a long sequence of isolated micro-steps.
\end{remark}

\begin{definition}[Reference-comparison optical signature]
\label{def:reference-comparison-optical-signature}
For a reduced retained-carrier model series $S_{\mathrm{mod}}(\lambda_i)$ and a same-grid reference series $S_{\mathrm{ref}}(\lambda_i)$, define the comparison signature
\[
\Sigma_{\mathrm{opt}}(\lambda_i):=S_{\mathrm{ref}}(\lambda_i)-S_{\mathrm{mod}}(\lambda_i).
\]
\end{definition}

\begin{proposition}[First data-anchored optical comparison principle]
\label{prop:first-data-anchored-optical-comparison-principle}
Assume the retained-carrier model has already passed the bundled first-order optical validation and that $S_{\mathrm{ref}}$ is used only as an external same-observable comparison anchor. Then $\Sigma_{\mathrm{opt}}$ already supports the first trichotomy: uniformly small discrepancy gives first-order compatibility with the chosen reference family, smooth systematic discrepancy gives a shifted/refined placement inside that family, and irregular or localized discrepancy must be passed back to structural diagnostics before any stronger identification claim is made.
\end{proposition}

\begin{proof}
This is the compressed operational reading of the validated retained-carrier optical package: once like is compared with like and the reference is not promoted to a second carrier, the discrepancy can only be read as compatibility, smooth family shift, or a diagnostic reopen signal.
\end{proof}

\begin{definition}[Reference-family comparison score]
\label{def:reference-family-comparison-score}
For a finite family $\{S_{\mathrm{ref}}^{(\alpha)}\}_{\alpha\in\mathcal A}$ on the same wavelength grid, let
\[
\mathfrak s(\alpha):=\mathfrak s\bigl(\Sigma^{(\alpha)}_{\mathrm{opt}}\bigr)
\]
be any first-order score that decreases when the discrepancy is small, smooth, and structurally stable.
\end{definition}

\begin{proposition}[First reference-family ranking principle]
\label{prop:first-reference-family-ranking-principle}
Under the same retained-carrier assumptions, the finite family $\mathcal A$ may already be ranked by $\mathfrak s(\alpha)$, with lower score interpreted as better first-order compatibility. The output is not a final identification theorem but an ordered shortlist of admissible candidate families.
\end{proposition}

\begin{proof}
The comparison principle yields a controlled discrepancy profile for each candidate family; a finite score extraction therefore produces an admissible first shortlist without changing the carrier-side framework.
\end{proof}

\begin{corollary}[First shortlist workflow for candidate optical families]
\label{cor:first-shortlist-workflow-for-candidate-optical-families}
The first optical workflow is: validate the reduced model series, compare it with each reference family, score the resulting discrepancies, and retain the lowest-scoring families as the first optical shortlist.
\end{corollary}

\begin{proof}
This is the direct operational form of Propositions~\ref{prop:first-data-anchored-optical-comparison-principle} and \ref{prop:first-reference-family-ranking-principle}.
\end{proof}

\begin{definition}[Optical shortlist back-reading pattern]
\label{def:optical-shortlist-back-reading-pattern}
A back-reading pattern is the first structural interpretation of the shortlist class, distinguishing smooth boundary-dominated, smooth carrier-dominated, smooth mixed boundary/carrier, and localized defect/readout-near mismatch patterns.
\end{definition}

\begin{proposition}[First structural back-reading principle from optical shortlists]
\label{prop:first-structural-back-reading-principle-from-optical-shortlists}
Every finite shortlist obtained from Corollary~\ref{cor:first-shortlist-workflow-for-candidate-optical-families} may already be read back into one of the first structural classes of Definition~\ref{def:optical-shortlist-back-reading-pattern}. Smooth uniform offset points to boundary reopening, smooth curvature drift points to carrier reopening, coexistence of both points to mixed reopening, and localized/sign-changing mismatch points to defect/readout reopening.
\end{proposition}

\begin{proof}
This is exactly the first inverse reading of the retained optical discrepancy: the geometric shape of the residual determines which structural sector must be reopened first.
\end{proof}

\begin{corollary}[First workflow from optical shortlist to structural diagnostic hint]
\label{cor:first-workflow-from-optical-shortlist-to-structural-diagnostic-hint}
The first shortlist-to-structure workflow is: build the shortlist, classify its residual shape, and pass the result to the corresponding structural reopening hint.
\end{corollary}

\begin{proof}
Immediate from Proposition~\ref{prop:first-structural-back-reading-principle-from-optical-shortlists}.
\end{proof}

\begin{definition}[Structural diagnostic reopening priority]
\label{def:structural-diagnostic-reopening-priority}
The reopening priority is the first finite ordering of admissible structural sectors induced by the back-reading class.
\end{definition}

\begin{proposition}[First diagnostic reopening-priority principle from optical back-reading]
\label{prop:first-diagnostic-reopening-priority-principle-from-optical-back-reading}
The first back-reading class determines the first reopening order: boundary first for smooth offset, carrier first for smooth curvature, mixed boundary/carrier first for smooth coexistence, and defect/readout first for localized or parity-sensitive triggers.
\end{proposition}

\begin{proof}
The back-reading class already singles out which sector best matches the observed discrepancy, so the coarse reopening order follows immediately.
\end{proof}

\begin{corollary}[First refinement workflow after optical back-reading]
\label{cor:first-refinement-workflow-after-optical-back-reading}
After optical back-reading, it is enough to open the leading sector first and postpone deeper alternative sectors unless the leading residual test fails.
\end{corollary}

\begin{proof}
Immediate from Definition~\ref{def:structural-diagnostic-reopening-priority} and Proposition~\ref{prop:first-diagnostic-reopening-priority-principle-from-optical-back-reading}.
\end{proof}

\begin{definition}[Structural bombe on a reopening menu]
\label{def:structural-bombe-on-a-reopening-menu}
Given a finite reopening menu $\mathcal M$, the structural bombe is the first elimination operator that removes candidates incompatible with the observed back-reading class, residual improvement test, and low-order admissibility constraints.
\end{definition}

\begin{proposition}[First structural bombe elimination principle on reopening menus]
\label{prop:first-structural-bombe-elimination-principle-on-reopening-menus}
On any finite reopening menu, the bombe eliminates every candidate that contradicts the active back-reading class or worsens the residual while preserving any candidate that stays structurally admissible and non-worsening.
\end{proposition}

\begin{proof}
This is the compressed elimination rule built into the first bombe layer: the menu shrinks by contradiction and non-improvement.
\end{proof}

\begin{definition}[Nested structural bombe]
\label{def:nested-structural-bombe}
A nested bombe is the same elimination mechanism applied to a finite submenu below a surviving parent branch.
\end{definition}

\begin{proposition}[Nested bombe backtracking principle]
\label{prop:nested-bombe-backtracking-principle}
If every child of a surviving parent branch is eliminated, then the parent branch itself must be reopened or discarded; otherwise the surviving child becomes the ordered next refinement path.
\end{proposition}

\begin{proof}
This is the finite backtracking rule for nested menus.
\end{proof}

\begin{definition}[First smooth mixed shortlist case]
\label{def:first-smooth-mixed-shortlist-case}
The first smooth mixed case is the shortlist regime where smooth offset and smooth curvature coexist, no leading localized defect trigger is visible, and the coarse menu contains the candidates $\mu_{\mathrm B},\mu_{\mathrm C},\mu_{\mathrm{BC}},\mu_{\mathrm D}$.
\end{definition}

\begin{proposition}[First bombe reduction in the smooth mixed optical case]
\label{prop:first-bombe-reduction-in-the-smooth-mixed-optical-case}
In the first smooth mixed case, the bombe eliminates $\mu_{\mathrm D}$ by absence of a defect trigger and eliminates the pure boundary and pure carrier branches whenever each leaves the complementary smooth residual visible. The surviving coarse branch is then $\mu_{\mathrm{BC}}$, and the first surviving child order is boundary correction before carrier-curvature refinement.
\end{proposition}

\begin{proof}
This is the first applied mixed-case elimination: smooth mixed evidence kills the defect branch and the two pure branches, leaving the ordered mixed refinement path.
\end{proof}

\begin{definition}[First defect/readout-near shortlist case]
\label{def:first-defect-readout-near-shortlist-case}
The first defect/readout-near case is the shortlist regime where the mismatch is localized or parity-sensitive and cannot be absorbed by smooth offset or smooth curvature alone.
\end{definition}

\begin{definition}[Readout-trigger sieve]
\label{def:readout-trigger-sieve}
The readout-trigger sieve is the first pre-bombe filter that retains only those candidates whose local trigger profile matches the observed localized defect/readout pattern.
\end{definition}

\begin{proposition}[First trigger-sieve-plus-bombe reduction in the defect/readout-near optical case]
\label{prop:first-trigger-sieve-plus-bombe-reduction-in-the-defect-readout-near-optical-case}
In the first defect/readout-near case, the trigger sieve removes the smooth candidates before the bombe stage, and the reduced bombe then retains the direct defect/readout branch whenever hybrid smooth-plus-readout branches fail to improve the localized trigger beyond the direct patch.
\end{proposition}

\begin{proof}
Localized trigger information is stronger than a smooth menu reading, so the sieve-first regime compresses the menu before the bombe makes the final coarse elimination.
\end{proof}

\begin{definition}[Minimal structural diagnostic transition tensor]
\label{def:minimal-structural-diagnostic-transition-tensor}
The minimal transition tensor is the four-axis array
\[
\Theta_{r,u,c}^{\;r'},
\]
with regime $r$, tool $u$, readout class $c$, and next regime $r'$, supported only on admissible structural moves.
\end{definition}

\begin{definition}[First diagnostic shell condition for the minimal transition tensor]
\label{def:first-diagnostic-shell-condition-for-the-minimal-transition-tensor}
A transition lies on the diagnostic shell if and only if it is simultaneously closure-compatible, transport-compatible, readout-compatible, and depth-compatible. Off-shell entries of $\Theta$ vanish.
\end{definition}

\begin{definition}[First explicit mini-assignment for the leading optical regimes]
\label{def:first-explicit-mini-assignment-for-the-leading-optical-regimes}
The first explicit mini-assignment sets the leading on-shell tool choice to bombe-first in the smooth mixed regime and sieve-first in the defect/readout-near regime.
\end{definition}

\begin{proposition}[Priority readout of the first explicit mini-assignment]
\label{prop:priority-readout-of-the-first-explicit-mini-assignment}
The mini-assignment already fixes the first local tool priority for the two leading regimes and therefore compresses the first branching node to a singleton tool choice.
\end{proposition}

\begin{proof}
This is the direct reading of the nonzero shell-supported entries of the mini-assignment.
\end{proof}

\begin{corollary}[What the explicit mini-assignment already buys for proof economy]
\label{cor:what-the-explicit-mini-assignment-already-buys-for-proof-economy}
Once the regime and readout class are identified, the first tool step may be cited directly from the mini-assignment instead of reopening the whole first tool menu.
\end{corollary}

\begin{proof}
Immediate from Proposition~\ref{prop:priority-readout-of-the-first-explicit-mini-assignment}.
\end{proof}

\begin{definition}[First explicit nested submenu assignment after the two leading entries]
\label{def:first-explicit-nested-submenu-assignment-after-the-two-leading-entries}
The first nested submenu assignment fixes the next ordered submenu move: boundary-offset repair before carrier-curvature repair in the smooth mixed line, and local defect/readout repair before any smooth supplement in the defect/readout-near line.
\end{definition}

\begin{proposition}[Second-step priority readout of the first nested submenu assignment]
\label{prop:second-step-priority-readout-of-the-first-nested-submenu-assignment}
The nested submenu assignment already compresses the second branching node to a singleton child priority on the two leading lines.
\end{proposition}

\begin{proof}
This is the second-step reading of the surviving nested submenu entries.
\end{proof}

\begin{definition}[First shell-supported branch-depth policy]
\label{def:first-shell-supported-branch-depth-policy}
The early branch-depth policy is the shell-supported rule that the first two branching depths are determined by the explicit menu assignment and the explicit nested submenu assignment.
\end{definition}

\begin{proposition}[Early-depth support propagation under the first branch-depth policy]
\label{prop:early-depth-support-propagation-under-the-first-branch-depth-policy}
Within the stabilized early window, support propagates exactly along the policy-selected branch and no competing sibling remains admissible.
\end{proposition}

\begin{proof}
This is the finite propagation rule induced by the first two singleton assignments.
\end{proof}

\begin{definition}[First shell-compatible local utility functional]
\label{def:first-shell-compatible-local-utility-functional}
Let $U$ be a local utility functional on on-shell candidates, combining residual gain, closure-compatibility, transport-compatibility, readout fidelity, and backtracking cost. Off-shell candidates are assigned utility $-\infty$.
\end{definition}

\begin{proposition}[Utility refinement of the early branch-depth policy]
\label{prop:utility-refinement-of-the-early-branch-depth-policy}
On the stabilized early window, the unique utility maximizer reproduces the existing singleton policy; beyond that window, utility provides the first ranked continuation rule on the remaining on-shell menu.
\end{proposition}

\begin{proof}
Within the early window the shell and the previous singleton policy already select one surviving move; utility therefore agrees there and extends the same ordering principle to later local menus.
\end{proof}

\begin{corollary}[What the first local utility layer already buys for later depth extension]
\label{cor:what-the-first-local-utility-layer-already-buys-for-later-depth-extension}
Later local continuation may be treated as shell-first, utility-second, and only afterwards as a hard singleton policy when stability persists.
\end{corollary}

\begin{proof}
Immediate from Proposition~\ref{prop:utility-refinement-of-the-early-branch-depth-policy}.
\end{proof}

\begin{definition}[First utility-guided provisional depth-$2$ continuation rule]
\label{def:first-utility-guided-provisional-depth-2-continuation-rule}
The provisional depth-$2$ rule ranks the admissible children of the first stabilized branch by the local utility functional.
\end{definition}

\begin{proposition}[Utility-supported continuation beyond the stabilized early-depth window]
\label{prop:utility-supported-continuation-beyond-the-stabilized-early-depth-window}
Beyond the first two stabilized depths, the next local move may be taken from the highest-ranked on-shell child so long as no break trigger, shell loss, or utility tie forces a reopening.
\end{proposition}

\begin{proof}
This is the first later-depth continuation principle induced by the utility layer.
\end{proof}

\begin{definition}[First explicit smooth-mixed depth-$2$ test case]
\label{def:first-explicit-smooth-mixed-depth-2-test-case}
In the smooth mixed line, the first depth-$2$ test case is the surviving mixed path after boundary-offset repair, with carrier-side curvature correction as the leading remaining local move.
\end{definition}

\begin{proposition}[First explicit third-step priority in the smooth mixed line]
\label{prop:first-explicit-third-step-priority-in-the-smooth-mixed-line}
The explicit smooth mixed depth-$2$ test case selects carrier-curvature follow-up as the unique third step.
\end{proposition}

\begin{proof}
This is the leading continuation left after the earlier mixed-case eliminations and utility ranking.
\end{proof}

\begin{definition}[First explicit defect/readout-near depth-$2$ test case]
\label{def:first-explicit-defectreadout-near-depth-2-test-case}
In the defect/readout-near line, the first depth-$2$ test case is the surviving direct readout path after the first local defect patch, with continued local defect/readout stabilization as the leading next move.
\end{definition}

\begin{proposition}[First explicit third-step priority in the defect/readout-near line]
\label{prop:first-explicit-third-step-priority-in-the-defectreadout-near-line}
The explicit defect/readout-near depth-$2$ test case selects continued local defect/readout stabilization as the unique third step.
\end{proposition}

\begin{proof}
This is the direct continuation of the sieve-first/readout-first line once the localized trigger remains the dominant residual feature.
\end{proof}

\begin{definition}[First QAM-coded local spectral menu operator]
\label{def:first-qam-coded-local-spectral-menu-operator}
For a finite shell-admissible local menu, define the diagonal QAM-coded menu operator whose diagonal entries are the corresponding local utility values.
\end{definition}

\begin{proposition}[Spectral singleton readout of shell-supported local continuation]
\label{prop:spectral-singleton-readout-of-shell-supported-local-continuation}
The finite spectral theorem for the local QAM-coded menu operator reproduces the utility winner as the top eigendirection whenever the local utility maximizer is unique.
\end{proposition}

\begin{proof}
A finite diagonal self-adjoint menu operator has the utility values as spectrum, so the unique maximizer is exactly the top local spectral readout.
\end{proof}

\begin{corollary}[What the first local spectral/QAM bridge already buys]
\label{cor:what-the-first-local-spectralqam-bridge-already-buys}
Local continuation can now be cited either as the unique utility winner or as the unique top eigendirection of the corresponding local QAM-coded menu operator.
\end{corollary}

\begin{proof}
Immediate from Proposition~\ref{prop:spectral-singleton-readout-of-shell-supported-local-continuation}.
\end{proof}
\begin{remark}[Status of the compressed later-depth controller surface in the present release line]
\label{rem:status-of-the-compressed-laterdepth-controller-surface-in-v324}
The detailed later-depth controller build-up of the preceding line can now be treated in compressed form on the active line. Instead of carrying the full micro-chain of propagation, break, reentry, decoder, action, consistency, patch, decomposition, concatenation/cut, and finite normal form as separate local interfaces, the present line packages the same finite admissible later-depth content into one exact citation surface.
\end{remark}

\begin{definition}[Compressed later-depth controller surface on admissible windows]
\label{def:compressed-laterdepth-controller-surface-on-admissible-windows}
Let
\[
I=[d_-,d_+]\cap\mathbb N
\]
be a finite admissible later-depth interval on the leading optical regimes. We say that $I$ is read on the \emph{compressed later-depth controller surface} if its later-depth control content is cited only through the token
\[
\mathfrak Z^{\mathrm{cyc}}(I)
=
\langle \mathsf{laterdepth},d_-,d_+,\mathfrak N^{\mathrm{cyc}}(I)\rangle,
\]
where $\mathfrak N^{\mathrm{cyc}}(I)$ is the finite controller normal form storing the ordered lengths of maximal propagated advance-patches together with the intervening separator symbols $\mathsf{reopen}$ and $\mathsf{anchor}$.

Equivalently, the compressed surface remembers exactly the finite alternation pattern
\[
\text{advance-patch}\; /\; \text{separator}\; /\; \text{advance-patch}\; /\; \cdots
\]
on $I$ and forgets only the inessential internal relabeling inside propagated patch interiors.
\end{definition}

\begin{proposition}[Exact finite control recovery from the compressed later-depth controller surface]
\label{prop:exact-finite-control-recovery-from-the-compressed-laterdepth-controller-surface}
Assume Definition~\ref{def:compressed-laterdepth-controller-surface-on-admissible-windows}. Let $I$ be a finite admissible later-depth interval.
\begin{enumerate}[leftmargin=2em,label=\textup{(\roman*)}]
  \item The token $\mathfrak Z^{\mathrm{cyc}}(I)$ determines the exact break-aware later-depth controller profile of $I$ up to internal depth-preserving relabeling inside maximal propagated advance-patches.
  \item Every later argument whose dependence on $I$ is only through that finite control profile may be written entirely on the compressed later-depth controller surface.
  \item Local gluing and cutting of adjacent admissible windows are recoverable on the compressed surface by concatenating and cutting the stored normal-form word at the corresponding boundary depth.
  \item Whenever a finite later-depth passage is used only through its controller pattern, citing $\mathfrak Z^{\mathrm{cyc}}(I)$ is exact and reopening the full micro-chain is unnecessary.
\end{enumerate}
\end{proposition}

\begin{proof}
The finite controller normal form already stores the unique alternating decomposition of $I$ into maximal propagated advance-patches and separator sites. The compressed token records exactly that normal form together with the interval endpoints, hence it carries the same finite controller content as the detailed later-depth micro-chain modulo the inessential relabeling inside propagated patch interiors. Therefore every argument depending only on this finite control pattern may be written against $\mathfrak Z^{\mathrm{cyc}}(I)$ alone. The local concatenation/cut behavior is recovered because gluing and cutting act directly on the same ordered normal-form word.
\end{proof}

\begin{corollary}[What the compressed later-depth controller surface already buys on the active line]
\label{cor:what-the-compressed-laterdepth-controller-surface-already-buys-on-the-active-line}
A substantial part of the later-depth controller apparatus is now available in one compact active citation layer. In particular, later arguments may cite a finite admissible later-depth window through
\[
\mathfrak Z^{\mathrm{cyc}}(I)
\]
instead of reopening the full sequence of propagation, break, reentry, decoder, action, consistency, controller, patch, decomposition, and concatenation/cut statements.
\end{corollary}

\begin{proof}
This is the immediate operational reading of Proposition~\ref{prop:exact-finite-control-recovery-from-the-compressed-laterdepth-controller-surface}.
\end{proof}

\begin{remark}[Status of the terminal interface-application citation token in the present v\docversion\ release line]
\label{rem:status-of-the-terminal-interface-application-surface-in-the-present-v324-line}
The present line still does \emph{not} claim one unrestricted final application theory covering every later use of the framework. It does, however, support one compressed terminal citation layer: whenever a later argument uses only the currently bundled interface-application package on one terminal crystal--shadow interface, that package may be cited through one terminal token instead of reopening the individual interface-level applications.
\end{remark}

\begin{definition}[Compressed terminal interface-application token]
\label{prop:terminal-citation-rule-for-the-present-v324-interface-application-line}
Let
\[
\mathfrak I_{\mathrm{v20}}
\bigl([\mathfrak C_1,\dots,\mathfrak C_n]_{\star_{\mathrm{cr}}\mathrm{CNF}}\bigr)
\]
be a terminal interface object on the present release line. We write
\[
\mathfrak T^{\mathrm{app}}(\mathfrak I)
=
\langle
\mathsf{terminal},
\mathfrak I,
\mathsf{defectface},
\mathsf{fluiddem},
\mathsf{carrierproj},
\mathsf{twoprobe}
\rangle
\]
if the later argument depends only on the bundled interface-application package of Proposition~\ref{prop:first-bundled-interface-application-package}, requires no still-deferred stronger claim, and uses no pre-interface data already forgotten by passage to the terminal surface.
\end{definition}

\begin{proposition}[Exact terminal citation on the compressed interface-application token]
\label{cor:operational-mini-checklist-for-the-terminal-interface-application-surface}
Assume Definition~\ref{prop:terminal-citation-rule-for-the-present-v324-interface-application-line}. If a later argument depends on a terminal interface object only through the four bundled interface-level applications encoded in \(\mathfrak T^{\mathrm{app}}(\mathfrak I)\), then citing \(\mathfrak T^{\mathrm{app}}(\mathfrak I)\) is exact and reopening the individual interface-level application principles is unnecessary. If any stronger generator, deeper compensator, explicit fluid coarse-graining claim, universal lensing claim, or forgotten pre-interface datum is needed, the proof must reopen the earliest stage where that information is still visible.
\end{proposition}

\begin{proof}
This is the bundled interface-application package of Proposition~\ref{prop:first-bundled-interface-application-package} rewritten as one terminal citation token. Exactness holds because the token stores precisely the four currently stabilized interface-level uses and nothing stronger; if a later argument needs more, it is by definition outside the token and must reopen the earlier stage carrying the missing information.
\end{proof}











\begin{remark}
This hypothesis does not identify the compensator literally with a
classical gravitational lens. Its weaker and structurally safer claim
is that the closure/HC core behaves as an effective focusing kernel for
probe traces --- the spectral-projector image $\mathrm{im}(C) =
\mathcal{S}_0$ of Theorem~\ref{thm:hc-spectral-projector} acts as the
common structural center around which both probe traces are organized.
\end{remark}

\subsection*{Why this is plausible in the present framework}

Three already established facts support the hypothesis.

First, the master equation admits distinguished structural projections to
the probe/Dirac, phase-shifted QFT, and sharpened Einstein regimes
(Proposition~\ref{prop:structural-lift-principle}). Hence the ART and QFT
probes are not foreign objects but two regime readings of one common
structural dynamics.

Second, the triadic two-probe system already contains an intrinsic third
channel $\Psi_\perp$ coupling the two probe modes. The middle sector
therefore carries genuine interaction data, not a mere bookkeeping slot.

Third, the Einstein reconstruction is already expressed through bilinears
of a closed triadic probe field, so probe traces are not external
diagnostics but the natural source of geometric data in the current
framework.

\subsection*{Trace observables}

\paragraph*{Relative deflection.}
Let $X^\mu_{\mathrm{ART}}(\tau)$ and $X^\mu_{\mathrm{QFT}}(\tau)$
denote effective probe traces. The deflection observable is
\[
  \Delta X^\mu(\tau)
  :=
  X^\mu_{\mathrm{ART}}(\tau) - X^\mu_{\mathrm{QFT}}(\tau).
\]

\paragraph*{Relative phase lag.}
\[
  \Delta\phi(\tau)
  :=
  \phi_{\mathrm{ART}}(\tau) - \phi_{\mathrm{QFT}}(\tau).
\]
This quantity is sensitive to compensator-mediated observer-frame
displacement (Remark~\ref{rem:frame-shift-spectral}).

\paragraph*{Cross-bilinear geometry channel.}
\[
  E^{a}_{\mu,\mathrm{cross}}
  :=
  \operatorname{Re}\bigl(
    \bar\Psi_{\mathrm{ART}}
    \,\Gamma^a\,
    \mathcal{D}_\mu
    \Psi_{\mathrm{QFT}}
  \bigr),
\]
and similarly with probes exchanged. These mixed quantities detect
geometry encoded not in a single probe sector but in the coupled
relation of both probes through the triadic core.

\paragraph*{Focusing functional.}
\[
  \mathcal{L}_{\mathrm{focus}}
  :=
  f\!\bigl(
    E^{a}_{\mu,\mathrm{cross}},\,
    \mathfrak{R}_C,\,
    \mathcal{H}_C
  \bigr),
\]
where $f$ detects concentration of the two traces toward the same
closure-stable branch. The expression is heuristic at this stage.

\subsection*{Expected regimes}

\begin{proposition}[Expected qualitative behaviour of probe lensing]
\label{prop:expected-probe-lensing-behaviour}
The bidirectional probe lensing hypothesis predicts:
\begin{enumerate}[label=(\roman*)]
  \item under exact closure, the trace structure sharpens and becomes
        sector-clean;
  \item when the triadic coupling channel is switched off, the
        cross-bilinear focusing signal weakens or disappears;
  \item in the phase-shifted observer regime, relative phase lag is more
        pronounced than geometric coframe data;
  \item in the large-scale closure-stable regime, relative focusing is
        expected to organize into an effective geometric readout.
\end{enumerate}
\end{proposition}

\begin{proof}
Exact closure suppresses leakage and makes the physical sector
autonomous (Lemma~\ref{lem:closure-forcing-lemma}), so trace observables
become sector-clean. The triadic two-probe equation shows that
$\Psi_\perp$ is the intrinsic coupling route; without it the two regime
probes largely decouple, weakening the cross-bilinear signal. By
Proposition~\ref{prop:structural-lift-principle}, QFT emerges in the
phase-shifted observer regime while the Einstein limit appears only at
large closure-stable scales, so phase effects dominate the QFT reading
and geometric focusing emerges only later.
\end{proof}

\subsection*{Relation to Einstein reconstruction}

The bidirectional probe picture may weaken the hardest current
bottleneck of the Einstein block. At present, the metric is reconstructed
from the bilinear coframe
\[
  E^a{}_\mu
  =
  \operatorname{Re}\bigl(
    \bar\Psi_\triangle\,\Gamma^a\,\mathcal{D}_\mu\Psi_\triangle
  \bigr),
\]
and the unresolved step is to justify its maximal rank from the triadic
probe dynamics itself. The bidirectional probe hypothesis suggests an
earlier readout layer: geometry may first appear as a stable trace
focusing/deflection pattern before it is compressed into a single
effective coframe.

\begin{remark}
If this route works, the Einstein sector would no longer require a
direct leap from closure to metric. The structural lift
(Definition~\ref{def:structural-lift}) of this two-step chain is:
\[
\begin{aligned}
\text{closure-stable probe traces}
&\Longrightarrow \text{effective focusing geometry}\\
&\Longrightarrow \text{coframe / metric reconstruction}.
\end{aligned}
\]
This could reduce the present model-dependence of the Einstein limit and
would realize the concrete form of the meta-field projection
(iv) in Hypothesis~\ref{hyp:hc-meta-generator}.
\end{remark}

% ------------------------------------------------------------------
% Theorem 7: HC-Induced Readout Geometry (v2.27.4 integration)
% ------------------------------------------------------------------

\section{HC-Induced Readout Geometry and Structural Unity}
\label{sec:theorem7}

\subsection*{Setup}

Let $\mathcal{K}$ be a closure-stable structural kernel, and let
\[
  \mathcal{R}_f : \mathcal{K} \to \mathcal{O}_f
\]
be admissible readout maps corresponding to different regimes (ART,
QFT, probe). Let $\mathcal{H}_C$ denote the Heisenberg compensator,
understood as a structural object controlling compatibility between
readouts. This is the concrete realization of
Definition~\ref{def:readout-relative-lemma}: each $\mathcal{R}_f$
produces a readout-relative statement, while $\mathcal{K}$ carries the
kernel-level content.

\subsection*{Theorem 7 (HC-Induced Readout Geometry)}

\begin{theorem}[HC-induced readout geometry]
\label{thm:hc-readout-geometry}
Under the above assumptions, the following statements hold.
\begin{enumerate}[label=(\arabic*)]
  \item \textbf{Readout divergence and effective trajectories.}
  For two readouts $f,g$, define
  $\Delta_{f,g} := \mathcal{R}_f - \mathcal{R}_g$.
  The induced effective dynamics in observable space generically
  produces trajectories that are not globally closed and may exhibit
  local segments resembling hyperbolic motion with apparent domain gaps.

  \item \textbf{Projection singularities.}
  Apparent gaps in trajectories correspond to incompatibilities of
  projections:
  \[
    \exists\, x \in \mathcal{K}:\quad
    \begin{aligned}
      &\mathcal{R}_f(x)\ \text{well-defined, while}\\
      &\mathcal{R}_g(x)\ \text{is not compatibly representable.}
    \end{aligned}
  \]
  These are projection singularities rather than physical
  discontinuities. In the language of
  Proposition~\ref{prop:legitimacy-restricted-lemmas}: $\mathcal{R}_f(x)$
  is a readout-relative statement whose regime of validity excludes the
  $g$-readout.

  \item \textbf{Mirror structure and recombination.}
  There exists a transformation $S$ such that locally
  $S(\mathcal{R}_f) \approx \mathcal{R}_g$.
  Complementary trajectory branches can therefore be recombined into
  an effective elliptic structure at the level of combined readouts.

  \item \textbf{Effective two-center dynamics.}
  The recombined dynamics is equivalent to motion in a symmetric
  two-center system, where the centers correspond to competing
  readout-fixed structures rather than physical sources. This is
  the bidirectional probe picture of
  Section~\ref{sec:bidirectional-probe-lensing}
  (Proposition~\ref{hyp:bidirectional-probe-lensing}) at the level of
  effective trajectories.

  \item \textbf{Non-orientable global structure.}
  The combined readout structure is generically non-orientable:
  transport along a closed loop induces a phase shift, requiring
  multiple loops for full return. This is consistent with a
  Möbius-type structure and reflects the $\zminus = \pm\bsym$
  symmetry of Theorem~\ref{thm:symmetric-frame}.

  \item \textbf{Existence of a structural lift.}
  There exists a lifted space $\tilde{\mathcal{K}}$ and projection
  $\pi : \tilde{\mathcal{K}} \to \mathcal{O}$ such that trajectories
  in $\tilde{\mathcal{K}}$ are smooth and closed, while their
  projections reproduce gaps, hyperbolic segments, and branching
  effects. This is the structural lift of
  Definition~\ref{def:structural-lift} applied to trajectory data.

  \item \textbf{HC as curvature.}
  The compensator $\mathcal{H}_C$ acts as curvature of the readout
  bundle, governing deviations of trajectories, phase shifts, and
  apparent failure of closure at the visible trajectory level. In the
  present release such failure must first be tested through native
  closure-class preservation and the OP transport theorem rather than
  inferred from raw trajectory separation alone. This is the geometric
  realization of Corollary~\ref{cor:M-HC-fixpoint}: $\mathcal{M}C = \lambda_0 C$
  fixes the compensator as the structural curvature anchor.

  \item \textbf{Projective vanishing.}
  There exist configurations such that
  \[
    \sum_i \mathcal{H}_C^{(f_i)} = 0
    \quad\text{while}\quad
    \mathcal{H}_C^{(\mathrm{kernel})} \neq 0.
  \]
  Thus the compensator may vanish in projection while remaining
  nontrivial at kernel level. This is precisely
  Remark~\ref{rem:hc-shadow-vs-core}: the shadow vanishes, not
  the kernel $C = \Pi^{(\mathcal{M})}_{\mathrm{spec}}$.

  \item \textbf{Emergent tensor structure.}
  The measurable curvature induced by $\mathcal{H}_C$ defines an
  effective tensor field $T_{\mu\nu}^{(\mathrm{HC})}$, arising from
  readout divergence $\Delta_{f,g}$ rather than from fundamental
  stress-energy. This is the tensorial compensator readout
  $\mathcal{K}_{\mu\nu}^{(\mathrm{HC})}$ of
  Hypothesis~\ref{hyp:hc-meta-generator}(iii) at the trajectory level.

  \item \textbf{Recursive field structure.}
  There exists a functional relation
  $\mathcal{F}[\mathcal{F}] = 0$,
  where $\mathcal{F}$ generates field equations recursively and
  $\mathcal{H}_C$ appears as a structural product within this
  recursion. This reflects the generating loop of
  Section~\ref{sec:hc-meta-field-generator}: the compensator is
  simultaneously generated by closure and inserted into transport.
\end{enumerate}
\end{theorem}

\begin{proof}
Statements (1)--(2) follow from the definition of $\Delta_{f,g}$ and
the readout-relative lemma framework
(Definition~\ref{def:readout-relative-lemma}): incompatible regime
projections produce apparent gaps that are projection singularities,
not structural discontinuities.

Statement (3) follows from the unitary equivalence of the three kernels
(Theorem~\ref{thm:unitary-equiv}): the transformation $S$ is the
unitary equivalence between the $f$- and $g$-readout representations.

Statement (4) is the trajectory-level reading of the bidirectional
probe picture (Section~\ref{sec:bidirectional-probe-lensing}): the two
regime probes approach the same closure-stable core from opposite sides.

Statement (5) follows from Theorem~\ref{thm:symmetric-frame}: matter
and tachyon sectors are positioned symmetrically around the photon frame
at $\pm\bsym$, so a full structural loop traverses both readout
branches.

Statement (6) is the structural lift principle
(Proposition~\ref{prop:structural-lift-principle}) applied to
trajectory data: the smooth lifted trajectories in $\tilde{\mathcal{K}}$
are the kernel-level content, while the projected gaps are readout
artifacts.

Statement (7) follows from Corollary~\ref{cor:M-HC-fixpoint}:
$\mathcal{M}C = \lambda_0 C$ identifies the compensator as the
structural fixpoint controlling all admissible transport; deviation
from this fixpoint is exactly what curvature measures.

Statement (8) is Remark~\ref{rem:hc-shadow-vs-core}: exact-sector
vanishing $\sum_i \mathcal{H}_C^{(f_i)} = 0$ is the shadow, not the
kernel $C = \Pi^{(\mathcal{M})}_{\mathrm{spec}}$
(Theorem~\ref{thm:hc-spectral-projector}).

Statement (9) follows from statement (7) and the probe-tensor readout
of Section~\ref{sec:bidirectional-probe-lensing}: the cross-bilinear
$E^a_{\mu,\mathrm{cross}}$ assembles into $T_{\mu\nu}^{(\mathrm{HC})}$.

Statement (10) is the recursive generating loop of
Hypothesis~\ref{hyp:hc-meta-generator}: the meta-field equation
$\mathfrak{E}_{\mathrm{HC}}[\ldots] = 0$ is the instance of
$\mathcal{F}[\mathcal{F}] = 0$ in the present framework.
\end{proof}

\subsection*{Interpretation}

Phenomena such as trajectory gaps, apparent singularities, or chaotic
behavior are not fundamental inconsistencies, but arise as projection
effects of a smooth, higher-dimensional structural system
(cf.\ Definition~\ref{def:kernel-level-statement}). Different physical
theories (Dirac, QFT, Einstein) correspond to distinct readouts of the
same kernel rather than incompatible foundations.

\subsection*{Methodological consequence}

Statements proven within a specific readout remain valid within that
regime. Their invariant structural content may be lifted to kernel-level
statements via structural lift, in accordance with
Section~\ref{sec:readout-relative-lemmas} and
Proposition~\ref{prop:structural-lift-principle}.

% ------------------------------------------------------------------
% Hypothesis 7.4 (v2.27.4 integration)
% ------------------------------------------------------------------
